\documentclass[11pt,a4paper]{article}
\usepackage[utf8]{inputenc}
% lmodern must precede fontenc: without a scalable Type1 font the T1 encoding
% falls back to bitmap fonts, and microtype's expansion then aborts the run
% with "auto expansion is only possible with scalable fonts".
\usepackage{lmodern}
\usepackage[T1]{fontenc}
\usepackage[margin=1in]{geometry}
\usepackage{amsmath,amssymb}
\usepackage{graphicx}
\usepackage{microtype}
\usepackage[hidelinks,breaklinks]{hyperref}
\usepackage{url}
\urlstyle{same}

% Generated by code/tools/md_to_tex.py from the SurveyForge Markdown output.
% Citation numbers match the source .md exactly.

\title{A Survey on the Optimization of Large Language Model-based Agents}
\author{Generated by SurveyForge}
\date{}

\begin{document}
\maketitle
\tableofcontents
\newpage



\section{Introduction}


The emergence of large language model-based agents has revolutionized the field of artificial intelligence, enabling machines to understand, generate, and interact with human language in unprecedented ways. As noted in \cite{arxiv:2206.07682}, the scaling up of language models has predictably improved performance and sample efficiency on a wide range of downstream tasks. This subsection provides an overview of the significance of large language model-based agents, their applications, and the importance of optimization in enhancing their performance and efficiency.

Large language models have been successfully applied in various domains, including natural language processing, decision-making, and problem-solving. For instance, \cite{doi:10.1016/j.metrad.2023.100017} presents a comprehensive survey of ChatGPT-related research, highlighting the potential of large language models in diverse applications. Similarly, \cite{doi:10.1145/3604237.3626869} explores the applications of large language models in finance, demonstrating their potential in tasks such as text classification, sentiment analysis, and machine translation.

The optimization of large language model-based agents is crucial in enhancing their performance and efficiency. As discussed in \cite{arxiv:2305.16264}, the scaling of language models in data-constrained regimes is an important area of research, highlighting the importance of optimizing model size and training data. Furthermore, \cite{arxiv:2309.03409} proposes a simple and effective approach to leverage large language models as optimizers, where the optimization task is described in natural language.

Despite the significant advancements in large language model-based agents, several challenges and limitations remain. For example, \cite{doi:10.1007/s10849-023-09409-x} evaluates the strengths and weaknesses of large language models, highlighting their potential biases and limitations. Similarly, \cite{doi:10.18653/v1/2023.acl-long.411} discusses the challenges and limitations of applying large language models to code-related tasks.

The future of large language model-based agents holds tremendous promise, with potential applications in emerging areas such as edge computing, human-computer interaction, and multimodal processing. As noted in \cite{arxiv:2401.13601}, multimodal large language models have undergone substantial advancements, enabling the integration of multiple modalities such as text, images, and audio. Furthermore, \cite{doi:10.1109/mnet.2024.3435752} explores the potential of large language models in networking applications, highlighting their potential in tasks such as network configuration and management.

In conclusion, large language model-based agents have revolutionized the field of artificial intelligence, enabling machines to understand, generate, and interact with human language in unprecedented ways. The optimization of these agents is crucial in enhancing their performance and efficiency, and various techniques such as fine-tuning, pruning, and knowledge distillation can be applied to achieve this goal. As research in this area continues to evolve, it is likely that large language model-based agents will play an increasingly important role in shaping the future of artificial intelligence.
The development of large language model-based agents has also been influenced by the concept of mixture of experts, as discussed in \cite{doi:10.1109/tkde.2025.3554028}. This approach has been shown to be effective in improving the performance of large language models, particularly in tasks that require specialized knowledge or expertise. Additionally, the use of large language models in multi-agent systems has been explored, as noted in \cite{doi:10.1109/mwc.016.2300600}, highlighting the potential of these models in enabling complex decision-making and problem-solving in dynamic environments.

The applications of large language model-based agents are diverse and widespread, ranging from natural language processing and decision-making to problem-solving and human-computer interaction. As discussed in \cite{arxiv:2309.14365}, the integration of large language models with multi-agent systems has the potential to enable more efficient and effective decision-making in complex networks. Furthermore, \cite{arxiv:2308.03688} provides a comprehensive overview of the evaluation of large language model-based agents, highlighting the importance of establishing comprehensive evaluation metrics for assessing their performance.

The optimization of large language model-based agents is a critical area of research, with significant potential for advancement and innovation. As noted in \cite{arxiv:2402.01364}, the development of continual learning techniques is essential in enabling large language models to adapt to changing environments and tasks. Furthermore, \cite{doi:10.1109/access.2024.3424945} provides a comprehensive overview of the techniques and strategies used to optimize the performance of large language models, highlighting the importance of model scaling, fine-tuning, and knowledge distillation.

Overall, the development of large language model-based agents has the potential to revolutionize the field of artificial intelligence, enabling machines to understand, generate, and interact with human language in unprecedented ways. The optimization of these agents is crucial in enhancing their performance and efficiency, and various techniques such as fine-tuning, pruning, and knowledge distillation can be applied to achieve this goal. As research in this area continues to evolve, it is likely that large language model-based agents will play an increasingly important role in shaping the future of artificial intelligence.


\section{Architectures and Designs for Large Language Model-based Agents}



\subsection{Introduction to Architectures for Large Language Model-based Agents}


The development of large language model-based agents has been a pivotal advancement in the field of artificial intelligence, enabling machines to understand, generate, and interact with human language in unprecedented ways. At the core of these agents are sophisticated architectures designed to process and generate human-like language, making them indispensable for tasks such as natural language processing, decision-making, and problem-solving. This subsection delves into the fundamental concepts and designs of these architectures, highlighting their significance and the challenges they pose.

Recent studies, such as \cite{arxiv:2206.07682}, have shown that scaling up language models can predictably improve performance and sample efficiency on a wide range of downstream tasks. However, this scaling also introduces unforeseen phenomena, such as emergent abilities that cannot be predicted by simply extrapolating the performance of smaller models. The existence of such emergence implies that additional scaling could further expand the range of capabilities of language models. Architectures like transformer-based models, which rely on self-attention mechanisms to process sequences of data, have been particularly successful in large language model-based agents. As discussed in \cite{doi:10.1016/j.metrad.2023.100017}, these models combine large-scale architectures with vast amounts of textual training data, allowing them to understand and generate text at a level that is often indistinguishable from human language.

The design of architectures for large language model-based agents involves several key considerations, including modularity, scalability, and adaptability. Modular designs, for instance, allow for the integration of different components, such as attention mechanisms and memory components, which can enhance the capabilities of the agent. Scalability is also crucial, as it enables the model to handle large amounts of data and compute resources efficiently. Adaptability is another important aspect, as it allows the agent to learn from feedback and adapt to new environments. As noted in \cite{doi:10.1145/3604237.3626869}, the choice of architecture can significantly impact the performance of the agent, and different architectures may be more suitable for different tasks and applications.

Despite the advancements in large language model-based agents, there are still several challenges that need to be addressed. One of the main challenges is the lack of interpretability and explainability of these models, making it difficult to understand how they arrive at their decisions. Another challenge is the potential for bias and unfairness in the models, which can perpetuate existing social inequalities. As discussed in \cite{doi:10.1145/3641289}, evaluating the performance of large language model-based agents is also a complex task, requiring careful consideration of metrics and benchmarks.

Recent studies have explored various approaches to addressing these challenges, including the use of techniques such as pruning, quantization, and knowledge distillation to reduce the computational costs of large language models. As noted in \cite{arxiv:2305.16264}, these techniques can help to improve the efficiency of the models without sacrificing performance. Other studies have focused on developing more interpretable and explainable models, such as \cite{arxiv:2309.02427}, which proposes a framework for designing cognitive architectures that can provide insights into the decision-making processes of large language model-based agents.

In conclusion, the design of architectures for large language model-based agents is a complex and multifaceted task that requires careful consideration of several key factors, including modularity, scalability, and adaptability. While there have been significant advancements in this field, there are still several challenges that need to be addressed, including the lack of interpretability and explainability, potential bias and unfairness, and the need for more efficient and effective evaluation metrics. As noted in \cite{doi:10.3390/app14052074}, the future of large language model-based agents holds much promise, but it will require continued innovation and research to address these challenges and realize the full potential of these models. Future studies, such as \cite{doi:10.34133/research.0646}, will be crucial in exploring the potential of large language models and addressing the challenges associated with their development and deployment.


\subsection{Key Components of Large Language Model-based Agent Architectures}


The development of large language model-based agents has been significantly enhanced by the incorporation of several key components, including attention mechanisms, memory components, and external knowledge integration. These components have revolutionized the capabilities of large language models, enabling them to generate coherent and contextually relevant text, and to incorporate external knowledge and generate more informative and accurate text. At the forefront of these components are attention mechanisms, which enable the model to focus on specific parts of the input sequence when generating text \cite{arxiv:1412.6980}. This is particularly evident in the transformer-based models, where self-attention mechanisms allow for the weighing of different input elements relative to each other. The effectiveness of attention mechanisms in improving model performance is well-documented, with studies demonstrating their ability to capture long-range dependencies and contextual relationships within input sequences \cite{arxiv:2402.15627}.

Building on the foundational concepts and designs of large language model-based agents discussed earlier, the integration of attention mechanisms and memory components has been shown to significantly enhance the performance of these models. Memory components, in particular, enable the model to store and retrieve information over extended periods, facilitating the generation of coherent and contextually relevant text \cite{doi:10.1145/3748302}. The integration of memory components into large language models has been demonstrated to improve their performance on a range of tasks, including question answering, text summarization, and dialogue generation. Furthermore, the use of external knowledge integration techniques, such as retrieval-augmented generation, has been shown to enhance the model's ability to incorporate external knowledge and generate more informative and accurate text \cite{arxiv:2401.02777}.

The incorporation of these key components into large language model-based agents has significant implications for their potential applications, particularly in areas such as natural language processing, human-agent collaboration, and decision-making. For instance, the development of agents capable of generating coherent and contextually relevant text has far-reaching implications for tasks such as language translation, text summarization, and dialogue generation. Moreover, the ability of these agents to incorporate external knowledge and generate informative text makes them particularly suited to tasks requiring the generation of specialized or technical content, such as scientific writing or educational materials. As the field continues to evolve, it is likely that we will see the development of even more sophisticated and capable large language model-based agents, with significant implications for a range of applications and industries.

However, despite the significant advancements that have been made in the development of large language model-based agents, several challenges and limitations remain. One of the primary concerns is the potential for these agents to generate biased or misleading text, which can have serious consequences in certain applications \cite{arxiv:2410.02644}. Additionally, the integration of external knowledge into these agents can be challenging, particularly when dealing with complex or nuanced topics. To address these challenges, researchers are exploring new techniques, such as multimodal learning and human-in-the-loop feedback, to enhance the performance and reliability of large language model-based agents. These challenges and limitations will be further explored in the following sections, which will discuss the optimization of large language model-based agents and the design of adaptive and modular architectures.

In conclusion, the key components of large language model-based agents, including attention mechanisms, memory components, and external knowledge integration, have significantly enhanced their capabilities and potential applications. As we move forward in the development of these agents, it is essential to address the challenges and limitations that remain, and to continue exploring new techniques and approaches to enhance their performance and reliability. The future of large language model-based agents holds much promise, and it will be exciting to see the advancements that are made in this field in the coming years, particularly in areas such as adaptive and modular architectures, which will be discussed in the following section.


\subsection{Designing Adaptive and Modular Architectures for Large Language Model-based Agents}


The development of large language model-based agents has led to a surge in research focused on designing adaptive and modular architectures that can efficiently integrate these models with other components, such as tools and knowledge graphs. As highlighted in \cite{arxiv:1906.09379}, the ability of large language models to capture complex patterns in natural language has significant implications for their use in agent architectures. However, the design of such architectures requires careful consideration of factors such as modularity, scalability, and adaptability.

One key approach to designing adaptive and modular architectures is the use of modular design principles, as discussed in \cite{arxiv:2308.04026}. This involves breaking down the architecture into smaller, independent components that can be easily integrated and modified. For example, \cite{doi:10.18653/v1/2023.acl-long.411} highlights the use of modular design principles in the development of large language models for coding tasks. By using modular design principles, developers can create architectures that are more flexible and adaptable to changing requirements.

Another important consideration in the design of adaptive and modular architectures is the use of attention mechanisms, as discussed in \cite{doi:10.18653/v1/2023.findings-emnlp.128}. Attention mechanisms allow the model to focus on specific parts of the input data, enabling more efficient processing and improved performance. \cite{arxiv:2301.04589} also highlights the importance of memory mechanisms in large language models, demonstrating how these mechanisms can be used to enhance the model's ability to process and generate text.

The integration of external knowledge and domain-specific information is also critical in the design of adaptive and modular architectures. As discussed in \cite{arxiv:2502.10708}, the use of domain-specific knowledge can significantly enhance the performance of large language models. For example, \cite{arxiv:2402.12914} demonstrates the use of domain-specific knowledge in human-agent collaboration tasks, highlighting the potential for large language models to enhance collaboration and decision-making.

In addition to these considerations, the design of adaptive and modular architectures must also take into account the potential risks and challenges associated with large language models. As highlighted in \cite{arxiv:2402.11208}, the use of large language models can introduce significant security risks, including backdoor threats and data poisoning. Therefore, developers must prioritize the design of secure and robust architectures that can mitigate these risks.

Emerging trends in the design of adaptive and modular architectures include the use of multimodal inputs, edge computing, and human-computer interaction. For example, \cite{doi:10.1109/mnet.2024.3435752} highlights the potential for large language models to enhance networking applications, including the use of multimodal inputs and edge computing. \cite{arxiv:2309.14365} also discusses the potential for large language models to enhance human-computer interaction, including the use of natural language processing and computer vision.

In conclusion, the design of adaptive and modular architectures for large language model-based agents is a complex and multifaceted challenge. By considering factors such as modularity, scalability, and adaptability, and by leveraging emerging trends and technologies, developers can create architectures that are more efficient, effective, and robust. As highlighted in \cite{arxiv:2406.04151}, the development of such architectures has significant implications for the future of artificial intelligence, including the potential for more generalizable and adaptable agents. Future research directions include the development of more secure and robust architectures, as well as the exploration of new applications and domains for large language model-based agents, such as those discussed in \cite{doi:10.22214/ijraset.2023.54677} and \cite{doi:10.1109/mwc.016.2300600}.


\subsection{Evaluating and Optimizing Architectures for Large Language Model-based Agents}


Evaluating and optimizing the performance of large language model-based agent architectures is a crucial step in ensuring their efficiency, effectiveness, and safety in various applications. As the development of large language model-based agents continues to evolve, the importance of evaluating and optimizing their performance cannot be overstated. The capabilities of large language models have been expanding rapidly, and their integration into agent architectures has opened up new avenues for autonomous decision-making and problem-solving, as discussed in \cite{arxiv:2206.07682}. However, this also introduces complexities in evaluation and optimization, as agents must balance competing objectives such as task completion, safety, and ethical considerations.

One of the primary challenges in evaluating large language model-based agents is the lack of standardized benchmarks and metrics. To address this challenge, \cite{arxiv:2308.03688} proposes a multi-dimensional benchmark framework to assess the performance of LLM-based agents across various tasks and environments. This work highlights the importance of comprehensive evaluation metrics that consider multiple facets of agent performance, including efficiency, effectiveness, and safety. Similarly, \cite{arxiv:2308.05960} introduces a framework for benchmarking and orchestrating LLM-augmented autonomous agents, emphasizing the need for scalable and flexible evaluation methodologies. These benchmarking frameworks are essential for comparing the performance of different agent architectures and identifying areas for improvement.

In addition to benchmarking, optimization techniques play a critical role in improving the performance of large language model-based agents. \cite{arxiv:2308.02151} presents a novel approach to optimizing LLM-based agents using policy gradient optimization. This method enables agents to learn from environment feedback and adapt to new tasks and environments. \cite{arxiv:2311.13884} also explores the use of actor-critic methods for optimizing LLM-based agents, demonstrating their potential in large-scale decision-making applications. These optimization techniques can significantly enhance the performance of large language model-based agents and enable them to operate effectively in complex environments.

Safety and ethical considerations are also critical aspects of evaluating and optimizing large language model-based agents. \cite{arxiv:2412.14470} introduces a comprehensive benchmark framework to evaluate the safety of LLM-based agents, highlighting the need for robust and reliable safety metrics. \cite{arxiv:2402.11208} also emphasizes the importance of addressing backdoor threats and ensuring the security of LLM-based agents. As large language model-based agents become increasingly autonomous, it is essential to ensure that they operate safely and ethically, and that their decision-making processes are transparent and explainable.

Emerging trends in evaluating and optimizing large language model-based agents include the use of multi-agent systems and graph-based approaches. \cite{arxiv:2501.05707} presents a novel approach to self-improvement in multi-agent systems, demonstrating the potential for diverse reasoning chains to enhance agent performance. \cite{arxiv:2402.16823} introduces a graph-based framework for optimizing LLM-based agents, highlighting the potential for scalable and efficient optimization methodologies. These emerging trends have significant implications for the future of large language model-based agents and highlight the need for continued research and development in this area.

In conclusion, evaluating and optimizing the performance of large language model-based agent architectures is a complex and multifaceted challenge. By leveraging advances in benchmarking, optimization techniques, and safety considerations, researchers can develop more efficient, effective, and safe agent architectures. As noted in \cite{doi:10.1145/3748302}, the development of more sophisticated memory mechanisms and evaluation metrics will be critical to advancing the field. Future research directions include the exploration of novel optimization techniques, the development of more comprehensive benchmark frameworks, and the investigation of safety and ethical considerations in large language model-based agents. By addressing these challenges and opportunities, researchers can unlock the full potential of large language model-based agents and drive progress in autonomous decision-making and problem-solving, ultimately leading to more sophisticated and effective agents that can operate in a wide range of scenarios and applications.


\subsection{Emerging Trends and Future Directions in Architectures for Large Language Model-based Agents}


The design of architectures for large language model-based agents is a rapidly evolving field, with emerging trends and future directions focused on enhancing the capabilities, efficiency, and interaction modes of these agents. One of the key trends is the integration of multimodal inputs, allowing agents to process and generate not just text but also images, audio, and other forms of data \cite{arxiv:2406.04692}. This multimodal approach is expected to significantly enhance the agents' ability to understand and interact with their environment, making them more versatile and effective in real-world applications.

Another significant trend is the application of edge computing to large language model-based agents \cite{doi:10.1109/mwc.016.2300600}. Edge computing enables the processing of data closer to where it is generated, reducing latency and improving real-time processing capabilities. This is particularly important for agents that need to respond quickly to changing conditions or user inputs. By leveraging edge computing, agents can achieve faster response times and more efficient data processing, making them more suitable for applications that require real-time interaction.

Human-computer interaction is also a critical area of focus for future architectures of large language model-based agents \cite{doi:10.1145/3748302}. As agents become more sophisticated, there is a growing need for intuitive and user-friendly interfaces that can effectively communicate with humans. This includes the development of natural language processing capabilities that can understand and respond to human language, as well as the creation of visual and auditory interfaces that can provide feedback and guidance to users.

In addition to these trends, there is a growing interest in the development of more transparent and explainable architectures for large language model-based agents \cite{doi:10.1609/aaaiss.v2i1.27691}. As agents become more autonomous and make decisions that can have significant consequences, there is a need for architectures that can provide insights into their decision-making processes. This includes the development of techniques for visualizing and interpreting the outputs of large language models, as well as the creation of frameworks that can provide explanations for the agents' actions and decisions.

The use of large language models as optimizers is another emerging trend in the field \cite{arxiv:2309.03409}. By leveraging the capabilities of large language models to generate and evaluate solutions to complex problems, agents can optimize their performance and achieve better outcomes. This includes the application of large language models to optimization problems in fields such as logistics, finance, and healthcare, where the ability to generate and evaluate solutions quickly and effectively can have a significant impact.

Furthermore, the development of agent-based systems that can learn and adapt in complex environments is a key area of research \cite{arxiv:2308.03688}. By creating agents that can learn from their interactions with the environment and adapt to changing conditions, researchers can develop more robust and effective systems that can operate in a wide range of scenarios. This includes the development of agents that can learn from feedback, adapt to new situations, and make decisions in real-time.

Finally, the integration of large language models with other AI technologies, such as computer vision and robotics, is expected to enable the development of more sophisticated and capable agents \cite{arxiv:2402.01680}. By combining the strengths of different AI technologies, researchers can create agents that can perceive and interact with their environment in a more human-like way, enabling a wide range of applications in fields such as healthcare, education, and customer service.

In conclusion, the design of architectures for large language model-based agents is a rapidly evolving field, with emerging trends and future directions focused on enhancing the capabilities, efficiency, and interaction modes of these agents. By leveraging advances in multimodal processing, edge computing, human-computer interaction, transparency, and explainability, and integration with other AI technologies, researchers can develop more sophisticated and effective agents that can operate in a wide range of scenarios and applications \cite{doi:10.1145/3641289}. As the field continues to evolve, it is likely that we will see the development of even more capable and versatile agents that can make a significant impact in a wide range of fields and industries \cite{arxiv:2403.14608}.


\subsection{Applications and Case Studies of Large Language Model-based Agent Architectures}


The applications and case studies of large language model-based agent architectures have been rapidly expanding, with a wide range of domains benefiting from their capabilities. Building on the technical advancements and emerging trends discussed earlier, such as the integration of multimodal inputs and the application of edge computing, large language model-based agents have been successfully applied to various natural language processing tasks, including text classification, sentiment analysis, and machine translation, as seen in \cite{arxiv:2402.06196}. Furthermore, the use of large language models as agents has also been explored in decision-making and problem-solving tasks, with studies like \cite{arxiv:2309.02427} demonstrating their potential in these areas. The integration of large language models with other technologies, such as computer vision and robotics, has enabled the development of more sophisticated agents that can interact with their environment in a more human-like way, as discussed in \cite{doi:10.18653/v1/2023.acl-long.411} and \cite{doi:10.1109/iccv51070.2023.00788}.

In addition to these applications, large language model-based agents have been used in various case studies to demonstrate their effectiveness in real-world scenarios, highlighting the need for more efficient and real-time processing capabilities, as well as the importance of human-computer interaction and transparency. For example, \cite{arxiv:2308.03688} presents a benchmark for evaluating the performance of large language models as agents, highlighting their strengths and weaknesses in different tasks. Similarly, \cite{arxiv:2308.05960} provides a comprehensive comparison of different large language model-based agents, showcasing their capabilities and limitations. These case studies not only demonstrate the potential of large language model-based agents but also identify challenges and areas for improvement, such as the need for more robust evaluation metrics and the development of more efficient training methods, as discussed in \cite{arxiv:2312.00678}.

The use of large language models as agents has also raised important questions about their potential impact on society, emphasizing the need for responsible and ethical development and evaluation practices. As discussed in \cite{doi:10.3348/kjr.2023.0997}, these models have the potential to revolutionize various industries, including healthcare, education, and finance. However, their development and deployment also raise concerns about bias, fairness, and transparency, as highlighted in \cite{doi:10.1016/j.tics.2023.08.006}. Therefore, it is essential to develop and evaluate large language model-based agents in a responsible and ethical manner, taking into account their potential consequences on individuals and society as a whole. Recent studies have also explored the use of large language models as agents in multi-agent systems, where they can interact with each other and with their environment to achieve complex tasks, as seen in \cite{arxiv:2310.20151}.

These systems have the potential to enable more efficient and effective collaboration between agents, leading to breakthroughs in areas such as robotics, autonomous vehicles, and smart cities. Moreover, the development of large language model-based agents has also led to the creation of new evaluation metrics and benchmarks, such as \cite{arxiv:2308.03688} and \cite{arxiv:2410.13825}, which can be used to assess their performance and identify areas for improvement. As the field continues to evolve, we can expect to see more sophisticated and capable large language model-based agents that can interact with their environment in a more human-like way, leading to breakthroughs in areas such as natural language processing, decision-making, and problem-solving, as discussed in \cite{arxiv:2402.12914} and \cite{doi:10.1145/3748302}. Ultimately, the future of large language model-based agents holds much promise, and their potential to revolutionize various aspects of our lives is vast, as highlighted in \cite{doi:10.1109/mnet.2025.3539338} and \cite{doi:10.1109/mnet.2024.3435752}.


\section{Optimization Techniques for Large Language Model-based Agents}



\subsection{Fine-Tuning Strategies for Large Language Model-based Agents}


Fine-tuning is a crucial step in optimizing the performance of large language model-based agents, as it allows them to learn task-specific patterns and relationships. This process involves adjusting the pre-trained model's parameters to fit the specific requirements of a task or environment. As noted in \cite{doi:10.1145/3641289}, fine-tuning can significantly improve the performance of large language models on downstream tasks. In the context of large language model-based agents, fine-tuning strategies can be broadly categorized into supervised fine-tuning and reinforcement learning from human feedback.

Supervised fine-tuning involves fine-tuning the pre-trained model on a specific task using labeled data. This approach has been shown to be effective in adapting large language models to various tasks, such as natural language processing and decision-making. For instance, \cite{doi:10.1145/3604237.3626869} demonstrates the potential of fine-tuning large language models for financial tasks, including text classification and sentiment analysis. However, supervised fine-tuning requires a significant amount of labeled data, which can be time-consuming and expensive to obtain. As highlighted in \cite{doi:10.1145/3649506}, the quality and quantity of training data play a crucial role in determining the performance of fine-tuned models.

Reinforcement learning from human feedback is another fine-tuning strategy that involves training the model using feedback from humans, such as rewards or penalties. This approach has been shown to be effective in improving the performance of large language models on tasks that require human-like understanding and decision-making. For example, \cite{arxiv:2402.10172} demonstrates the potential of reinforcement learning from human feedback in optimizing large language models for complex tasks. However, this approach requires a significant amount of human feedback, which can be time-consuming and expensive to obtain. As noted in \cite{arxiv:2409.06857}, the use of small models can help reduce the computational costs associated with reinforcement learning from human feedback.

Few-shot learning is another fine-tuning strategy that involves fine-tuning the model on a limited number of examples. This approach has been shown to be effective in adapting large language models to new tasks with limited training data. For instance, \cite{arxiv:2206.07682} demonstrates the potential of few-shot learning in improving the performance of large language models on tasks that require human-like understanding and decision-making. However, few-shot learning requires a significant amount of prior knowledge and can be sensitive to the quality of the training data. As highlighted in \cite{doi:10.3390/app14052074}, the use of few-shot learning can help improve the performance of large language models on tasks that require human-like understanding and decision-making.

In recent years, there has been a growing interest in developing more efficient fine-tuning strategies for large language models. For example, \cite{arxiv:2402.05120} demonstrates the potential of using multiple agents to fine-tune large language models. This approach has been shown to be effective in improving the performance of large language models on tasks that require human-like understanding and decision-making. As noted in \cite{arxiv:2406.04692}, the use of multiple agents can help improve the performance of large language models on tasks that require human-like understanding and decision-making.

In conclusion, fine-tuning is a crucial step in optimizing the performance of large language model-based agents. Supervised fine-tuning, reinforcement learning from human feedback, and few-shot learning are some of the fine-tuning strategies that can be used to adapt pre-trained models to specific tasks and environments. While these strategies have been shown to be effective, they require significant amounts of labeled data, human feedback, or prior knowledge. Therefore, developing more efficient fine-tuning strategies that can adapt pre-trained models to new tasks with limited training data is an active area of research. As highlighted in \cite{doi:10.34133/research.0655}, the development of more efficient fine-tuning strategies can help improve the performance of large language models on tasks that require human-like understanding and decision-making. Future research directions include developing more efficient fine-tuning strategies, such as using small models or multiple agents, and exploring the potential of few-shot learning and reinforcement learning from human feedback. As noted in \cite{arxiv:2404.01023}, the use of multiple agents can help improve the performance of large language models on tasks that require human-like understanding and decision-making.


\subsection{Model Compression Techniques for Large Language Model-based Agents}


Model compression techniques have become essential for deploying large language model-based agents in resource-constrained environments, such as mobile devices or embedded systems. These techniques aim to reduce the computational costs and memory requirements of large language models without sacrificing their performance. One of the most popular model compression techniques is pruning, which involves removing redundant or unnecessary weights and connections in the model \cite{doi:10.1162/coli_a_00107}. This approach has been shown to be effective in reducing the size of large language models while preserving their accuracy.

In addition to pruning, other model compression techniques have been proposed, including quantization, which involves reducing the precision of the model's weights and activations \cite{arxiv:1412.6980}. Quantization can be applied to both the model's weights and activations, resulting in significant reductions in memory usage and computational costs. Another model compression technique is knowledge distillation, which involves transferring knowledge from a large pre-trained model to a smaller model \cite{arxiv:2402.15627}. This approach has been shown to be effective in reducing the size of large language models while preserving their performance.

Other model compression methods have also been proposed, including low-rank approximation, which involves approximating the weight matrix of a model using a low-rank matrix \cite{arxiv:2402.06196}, and sparse coding, which involves representing the weights of a model using a sparse code \cite{doi:10.3390/app14052074}. These techniques can be used individually or in combination to achieve optimal results. The choice of model compression technique depends on the specific requirements of the application and the characteristics of the model.

Recent studies have demonstrated the effectiveness of model compression techniques in reducing the computational costs and memory requirements of large language models. For example, \cite{arxiv:2112.11446} demonstrated that pruning and quantization can reduce the computational costs of large language models by up to 90\% without sacrificing their performance. \cite{doi:10.18653/v1/2022.bigscience-1.1} demonstrated that knowledge distillation can reduce the size of large language models by up to 90\% while preserving their performance. These results highlight the importance of model compression techniques in enabling the deployment of large language model-based agents in resource-constrained environments.

In conclusion, model compression techniques are essential for deploying large language model-based agents in resource-constrained environments. Various techniques, including pruning, quantization, and knowledge distillation, can be used to reduce the computational costs and memory requirements of large language models without sacrificing their performance. The choice of technique depends on the specific requirements of the application and the characteristics of the model. Further research is needed to develop more effective model compression techniques and to apply these techniques to a wider range of applications \cite{arxiv:2401.02777}. As the field continues to evolve, we can expect to see the development of more sophisticated model compression techniques and the application of these techniques to a wider range of domains, including natural language processing, computer vision, and robotics \cite{doi:10.1145/3748302}.


\subsection{Multi-Task Learning and Transfer Learning for Large Language Model-based Agents}


The use of multi-task learning, transfer learning, and meta-learning has emerged as a crucial aspect of optimizing large language model-based agents. These techniques enable the models to learn from multiple tasks, adapt to new environments, and improve their overall performance and efficiency. As discussed in \cite{arxiv:2309.14365}, the ability of large language models to generalize across tasks and domains is essential for their application in real-world scenarios.

Multi-task learning involves training a single model on multiple tasks simultaneously, allowing it to learn shared representations and improve its performance on each task. This approach has been shown to be effective in \cite{arxiv:1906.09379} and \cite{doi:10.18653/v1/2021.naacl-main.341}, where multi-task learning was used to improve the performance of large language models on various natural language processing tasks. For instance, \cite{doi:10.18653/v1/2023.acl-long.411} demonstrated that multi-task learning can be used to improve the performance of large language models on code generation tasks.

Transfer learning, on the other hand, involves pre-training a model on a large dataset and then fine-tuning it on a smaller dataset specific to the target task. This approach has been widely used in \cite{arxiv:2304.05332} and \cite{arxiv:2309.02427}, where transfer learning was used to adapt large language models to new tasks and domains. For example, \cite{arxiv:2308.03688} used transfer learning to evaluate the performance of large language models on various agent-based tasks.

Meta-learning involves training a model to learn how to learn from new tasks, allowing it to adapt quickly to new environments and tasks. This approach has been explored in \cite{arxiv:2205.05718} and \cite{arxiv:2301.04589}, where meta-learning was used to improve the performance of large language models on out-of-distribution tasks.

The use of multi-task learning, transfer learning, and meta-learning has several benefits, including improved performance, increased efficiency, and enhanced adaptability. However, these approaches also have their limitations and challenges, such as the need for large amounts of training data, the risk of overfitting, and the difficulty of selecting the most effective tasks and datasets for training.

Despite these challenges, the use of multi-task learning, transfer learning, and meta-learning has the potential to revolutionize the field of large language model-based agents. As discussed in \cite{arxiv:2306.03917} and \cite{arxiv:2406.04692}, these techniques can be used to create more generalizable, adaptable, and human-like language models.

In conclusion, the use of multi-task learning, transfer learning, and meta-learning is a crucial aspect of optimizing large language model-based agents. These techniques enable the models to learn from multiple tasks, adapt to new environments, and improve their overall performance and efficiency. While there are challenges and limitations to these approaches, they have the potential to revolutionize the field of large language model-based agents and enable the creation of more generalizable, adaptable, and human-like language models, as seen in \cite{arxiv:2402.12914} and \cite{arxiv:2407.07086}. Future research should focus on exploring the potential of these techniques and addressing the challenges and limitations associated with their use, as discussed in \cite{arxiv:2303.18223} and \cite{arxiv:2502.10708}.


\subsection{Optimization Algorithms for Large Language Model-based Agents}


The optimization of large language model-based agents is a crucial aspect of their development, as it directly impacts their ability to learn from data and adapt to new environments. Building on the foundations established in the previous sections, including model compression techniques and multi-task learning, the choice of optimization algorithm plays a vital role in the optimization process. In this context, various algorithms have been explored, including stochastic gradient descent, Adam, and Adagrad, as discussed in \cite{arxiv:2206.07682}. The selection of an appropriate optimization algorithm is vital, as it affects the model's convergence rate, stability, and overall performance.

Stochastic gradient descent (SGD) is a widely used optimization algorithm in the context of large language models. It involves updating the model's parameters using the gradient of the loss function, computed from a random subset of the training data. SGD is simple to implement and has been shown to be effective in optimizing large language models, as demonstrated in \cite{arxiv:2308.02151}. However, SGD can be sensitive to the choice of hyperparameters, such as the learning rate, and may require careful tuning to achieve optimal results. Furthermore, the use of SGD in conjunction with multi-task learning and transfer learning can lead to improved performance and efficiency, as discussed in \cite{doi:10.18653/v1/2023.acl-long.411} and \cite{arxiv:2308.03688}.

In contrast, Adam and Adagrad are adaptive optimization algorithms that adjust the learning rate for each parameter based on the magnitude of the gradient. Adam, in particular, has been widely adopted in the optimization of large language models, due to its ability to adapt to the scale of the gradients, as discussed in \cite{arxiv:2311.13884}. Adagrad, on the other hand, has been shown to be effective in optimizing models with sparse gradients, as demonstrated in \cite{arxiv:2311.06622}. Recent studies have also explored the use of more advanced optimization algorithms, such as momentum-based methods and quasi-Newton methods, as discussed in \cite{arxiv:2406.04151}. These algorithms have been shown to improve the convergence rate and stability of large language models, particularly in the context of multi-agent systems.

In addition to the choice of optimization algorithm, other factors can also impact the optimization of large language model-based agents. For example, the use of regularization techniques, such as dropout and weight decay, can help to prevent overfitting and improve the model's generalization performance, as discussed in \cite{doi:10.1609/aaaiss.v2i1.27688}. The selection of hyperparameters, such as the batch size and learning rate, can also significantly impact the optimization process, as demonstrated in \cite{arxiv:2401.06201}. Moreover, the evaluation of optimization algorithms for large language model-based agents is also an important consideration, as traditional evaluation metrics may not be sufficient to capture the complexities of these models, as discussed in \cite{arxiv:2412.14470}.

The evaluation of optimization algorithms is closely tied to the overall performance of large language model-based agents, which will be discussed in the following section on evaluation metrics. The choice of optimization algorithm and evaluation metric can significantly impact the model's performance, and the integration of human feedback and evaluation has also been explored as a means to improve the model's performance and adaptability, particularly in applications where human-agent collaboration is critical \cite{arxiv:2402.12914}. By exploring the use of advanced optimization algorithms and techniques, such as those discussed in \cite{arxiv:2410.06153}, researchers can improve the performance and efficiency of large language model-based agents, enabling them to be deployed in a wide range of applications, from natural language processing to decision-making and problem-solving, as discussed in \cite{doi:10.1145/3748302}. As the field continues to evolve, it is likely that new optimization algorithms and techniques will be developed, further improving the capabilities of large language model-based agents, as discussed in \cite{arxiv:2401.10910}.


\subsection{Evaluation Metrics for Large Language Model-based Agents}


Evaluating the performance of large language model-based agents is a multifaceted challenge that requires careful consideration of various metrics. The choice of evaluation metric is crucial, as it directly impacts the model's ability to learn from data and adapt to new environments. Perplexity, accuracy, and F1-score are commonly used metrics to assess the performance of large language model-based agents \cite{doi:10.1145/3641289}. Perplexity measures the model's ability to predict the next word in a sequence, while accuracy and F1-score evaluate the model's performance on specific tasks such as text classification and question answering \cite{arxiv:2404.09135}.

When evaluating large language model-based agents, it is essential to consider the trade-offs between different metrics. For instance, optimizing for perplexity may not necessarily lead to better performance on downstream tasks \cite{doi:10.1145/3748302}. Moreover, the choice of evaluation metric may depend on the specific application or task, highlighting the need for a nuanced understanding of the strengths and limitations of each metric \cite{arxiv:2402.12914}. Recent studies have also explored the use of alternative metrics, such as those based on human evaluations or behavioral analysis, to provide a more comprehensive understanding of agent performance \cite{arxiv:2308.03688}.

The development of evaluation metrics for large language model-based agents is an active area of research, with emerging trends focusing on the creation of more nuanced and task-specific metrics \cite{arxiv:2404.09022}. For example, the use of adversarial testing and robustness metrics has been proposed to evaluate the model's ability to withstand attacks or perturbations \cite{arxiv:2402.11208}. Additionally, the integration of human feedback and evaluation has been explored as a means to improve the model's performance and adaptability \cite{arxiv:2402.12914}.

Despite the progress made in developing evaluation metrics for large language model-based agents, several challenges and limitations remain. One of the primary concerns is the lack of standardization and consistency in evaluation protocols, which can make it difficult to compare the performance of different models or agents \cite{doi:10.1145/3641289}. Furthermore, the increasing complexity and scale of large language models pose significant challenges for evaluation, highlighting the need for more efficient and effective methods \cite{arxiv:2312.00678}.

To address these challenges, researchers have proposed various solutions, including the development of more efficient evaluation protocols, the use of surrogate models or proxies, and the integration of human evaluation and feedback \cite{arxiv:2410.06153}. Additionally, the creation of benchmark datasets and evaluation frameworks, such as AgentBench and WebArena, has facilitated the comparison and evaluation of large language model-based agents \cite{arxiv:2308.03688}. These advancements have the potential to significantly improve the evaluation and optimization of large language model-based agents, enabling more effective and efficient development of these models.

In conclusion, the evaluation of large language model-based agents is a complex and multifaceted challenge that requires careful consideration of various metrics and approaches. By understanding the strengths and limitations of different evaluation metrics and protocols, researchers can develop more effective and efficient methods for optimizing and evaluating the performance of these models. As the field continues to evolve, it is essential to address the challenges and limitations associated with evaluation, including the lack of standardization, consistency, and efficiency. By doing so, researchers can unlock the full potential of large language model-based agents and enable their widespread adoption in a variety of applications \cite{arxiv:2303.18223}. Future research directions may include the development of more nuanced and task-specific evaluation metrics, the integration of human feedback and evaluation, and the creation of more efficient and effective evaluation protocols \cite{arxiv:2309.03409}.


\subsection{Challenges and Future Directions for Large Language Model-based Agents}


The development and deployment of large language model-based agents have made significant strides in recent years, with applications in natural language processing, decision-making, and problem-solving \cite{arxiv:2402.06196}. As these agents continue to evolve, it is essential to address the challenges that must be overcome to optimize their performance and ensure their safe and effective deployment in real-world applications. One of the primary challenges is scalability, as large language models require substantial computational resources and memory to operate efficiently \cite{arxiv:2312.00678}. This challenge is closely tied to the evaluation metrics discussed in the previous section, as the choice of metric can significantly impact the model's performance and scalability.

Another significant challenge is interpretability, as large language models are often considered black boxes, making it difficult to understand their decision-making processes \cite{doi:10.1016/j.tics.2023.08.006}. To address this issue, researchers have proposed various techniques, including attention mechanisms, feature importance, and model explainability methods, which can provide insights into the models' internal workings \cite{arxiv:2309.02427}. Furthermore, the development of more transparent and explainable models, such as multimodal large language models, can also improve interpretability \cite{arxiv:2401.13601}. The importance of interpretability is highlighted by the need for more nuanced and task-specific evaluation metrics, as discussed in the previous section, which can help to uncover the strengths and limitations of large language model-based agents.

Robustness is another critical challenge, as large language models can be vulnerable to adversarial attacks, data poisoning, and other forms of manipulation \cite{doi:10.1145/3637528.3671650}. To address this issue, researchers have proposed various techniques, including adversarial training, data augmentation, and robustness evaluation methods, which can improve the models' resilience to attacks and errors \cite{doi:10.18653/v1/2023.acl-long.411}. Additionally, the development of more robust and reliable models, such as those using ensemble methods or uncertainty estimation, can also improve their overall performance \cite{doi:10.1145/3748302}. The robustness of large language models is closely tied to their evaluation, as the choice of evaluation metric can significantly impact the model's ability to withstand attacks or perturbations.

In terms of future directions, researchers are exploring various avenues to improve the performance and efficiency of large language model-based agents. One promising area is the development of multimodal large language models, which can integrate multiple sources of information, such as text, images, and audio, to improve their understanding and generation capabilities \cite{doi:10.1109/bigdata59044.2023.10386743}. Another area is the use of reinforcement learning and other optimization techniques to fine-tune large language models for specific tasks and environments \cite{doi:10.1109/tnnls.2024.3497992}. Furthermore, the development of more efficient and scalable models, such as those using sparse attention or parallel processing, can also improve their deployment in real-world applications \cite{doi:10.1145/3754448}. These advancements will be crucial in enabling the widespread adoption of large language model-based agents in various applications, including those discussed in the following section.

The integration of large language models with other AI technologies, such as computer vision and robotics, is also a promising area of research \cite{doi:10.1109/mnet.2024.3435752}. This can enable the development of more sophisticated and autonomous agents that can interact with their environment and perform complex tasks \cite{arxiv:2402.12914}. Additionally, the use of large language models in edge computing and other distributed computing paradigms can also improve their efficiency and scalability \cite{doi:10.1109/mnet.2025.3539338}. As researchers continue to advance the field, it is crucial to prioritize transparency, explainability, and robustness to ensure that large language model-based agents can be trusted and relied upon in critical applications \cite{arxiv:2404.01023}. By addressing the challenges associated with their development and deployment, large language model-based agents can revolutionize various aspects of our lives, from natural language processing to decision-making and problem-solving, and enable new applications and use cases that are yet to be explored.


\section{Evaluation Metrics and Benchmarks for Large Language Model-based Agents}



\subsection{Introduction to Evaluation Metrics for Large Language Model-based Agents}


The development and deployment of large language model-based agents have ushered in a new era of artificial intelligence, with these models demonstrating unprecedented capabilities in natural language understanding and generation. As these agents become increasingly pervasive, the need for comprehensive evaluation metrics to assess their performance has become paramount. Evaluation metrics play a crucial role in determining the efficacy of large language model-based agents, enabling researchers and practitioners to identify areas of improvement, compare different models, and ensure that these agents operate within desired parameters \cite{doi:10.1145/3641289}. 

The choice of evaluation metrics is multifaceted, depending on the specific application, task requirements, and desired outcomes. Common evaluation metrics include perplexity, accuracy, F1-score, and success rate, among others. However, these metrics have their limitations, and there is a growing recognition of the need for more nuanced and context-dependent evaluation frameworks \cite{arxiv:2402.06196}. For instance, \cite{arxiv:2206.07682} highlights the importance of evaluating large language models based on their emergent abilities, which are not present in smaller models but emerge as these models scale up.

Recent studies have also emphasized the role of benchmarks in evaluating the performance of large language model-based agents. Benchmarks such as AgentBench provide standardized frameworks for assessing agent capabilities across various tasks and environments \cite{arxiv:2308.03688}. These benchmarks not only facilitate comparisons between different models but also help identify areas where large language model-based agents excel or struggle, thereby guiding future research and development efforts.

Moreover, the evaluation of large language model-based agents extends beyond traditional metrics and benchmarks. As these agents interact with humans and other systems, aspects such as safety, ethics, and transparency become critical \cite{doi:10.1145/3604237.3626869}. Ensuring that these agents operate in a manner that is fair, accountable, and respectful of human values is essential for their widespread adoption and trustworthiness \cite{doi:10.1109/access.2024.3365742}.

Innovative approaches to evaluation are also being explored, including the use of large language models themselves as tools for evaluating other models \cite{arxiv:2212.09251}. This self-referential approach to evaluation highlights the complex and evolving nature of large language model-based agents and the need for adaptive and multifaceted evaluation strategies.

As the field continues to evolve, there is a pressing need for ongoing research into evaluation metrics and benchmarks that can keep pace with the rapid advancements in large language model technology \cite{arxiv:2402.01364}. This includes exploring new evaluation tasks, developing more sophisticated metrics that capture the nuances of agent performance, and addressing the challenges associated with evaluating agents in dynamic and real-world environments.

In conclusion, the evaluation of large language model-based agents is a multifaceted challenge that requires a comprehensive and dynamic approach. By leveraging a range of evaluation metrics, benchmarks, and innovative strategies, researchers and practitioners can ensure that these agents meet the highest standards of performance, safety, and ethical responsibility. As we move forward in this rapidly evolving field, the development of robust and adaptable evaluation frameworks will be crucial for unlocking the full potential of large language model-based agents and realizing their benefits across various domains and applications. Future research directions should focus on addressing the current limitations of evaluation metrics, exploring new benchmarks and evaluation tasks, and developing more sophisticated methods for assessing the complex behaviors of large language model-based agents.


\subsection{Benchmarking Large Language Model-based Agents}


Benchmarking large language model-based agents is a critical step in evaluating their performance, identifying areas for improvement, and comparing different models. As the field of large language model-based agents continues to evolve, the development of standardized benchmarks and evaluation protocols is essential for ensuring the fairness and reliability of these comparisons. The lack of standardized benchmarks has hindered the progress of large language model-based agents, as noted in \cite{arxiv:2309.14365}. However, recent advances in benchmarking and evaluation protocols have paved the way for the development of more effective and efficient large language model-based agents.

Recently, benchmarks such as AgentBench \cite{arxiv:2308.03688} and ShortcutsBench have been proposed to address the issue of standardized benchmarks. These benchmarks provide a comprehensive evaluation framework for assessing the performance of large language model-based agents in various tasks, including reasoning, problem-solving, and decision-making. For instance, AgentBench consists of eight distinct environments that test the agent's ability to reason, solve problems, and make decisions in complex scenarios. The benchmark also provides a set of evaluation metrics, including success rate, completion time, and efficiency, to measure the agent's performance.

The importance of benchmarking large language model-based agents is further emphasized in \cite{arxiv:2402.14744}, which highlights the need for standardized evaluation protocols to compare the performance of different agents. The study proposes a framework for evaluating the performance of large language model-based agents in urban mobility tasks, including route planning and navigation. Additionally, \cite{arxiv:2402.16142} provides a comprehensive review of large language models' versatility in various tasks, including text classification, sentiment analysis, and machine translation.

Furthermore, \cite{doi:10.3390/app14052074} provides a comprehensive survey of current trends, techniques, and challenges in large language models, including benchmarking and evaluation protocols. \cite{arxiv:2402.01680} also notes that the lack of standardized benchmarks has hindered the progress of large language model-based agents. However, the development of standardized benchmarks and evaluation protocols is an active area of research, with recent studies proposing various benchmarks and evaluation protocols for large language model-based agents.

In conclusion, benchmarking large language model-based agents is crucial for evaluating their performance, identifying areas for improvement, and comparing different models. The development of standardized benchmarks and evaluation protocols is essential for ensuring the fairness and reliability of these comparisons. As the field continues to evolve, it is likely that new benchmarks and evaluation protocols will be developed to address the emerging challenges and opportunities in large language model-based agents, ultimately leading to more effective and efficient agents that can be deployed in a wide range of applications, as discussed in the following section.


\subsection{Evaluation Metrics for Agent Performance}


The evaluation of large language model-based agents necessitates a comprehensive framework that assesses their performance across multiple dimensions, including task completion, safety, and ethical considerations. As highlighted in \cite{arxiv:1906.09379}, the development of evaluation metrics for these agents is crucial, given their potential to revolutionize various aspects of human-computer interaction. A key challenge in this context is the design of metrics that can effectively capture the nuances of agent performance, particularly in complex, dynamic environments.

One approach to evaluating large language model-based agents is through the use of benchmarks, such as AgentBench \cite{arxiv:2308.03688} and WebArena, which provide a standardized framework for assessing agent capabilities. These benchmarks typically involve a series of tasks that test the agent's ability to reason, problem-solve, and interact with its environment. For instance, \cite{doi:10.18653/v1/2021.naacl-main.341} demonstrates the potential of large language models in generating coherent, context-dependent text, which is a critical aspect of agent performance.

In addition to benchmarks, various evaluation metrics have been proposed to assess the performance of large language model-based agents. These metrics can be broadly categorized into three types: task-oriented, safety-oriented, and ethics-oriented. Task-oriented metrics, such as success rate and completion time, focus on the agent's ability to complete specific tasks \cite{doi:10.1038/s41746-024-01083-y}. Safety-oriented metrics, on the other hand, aim to evaluate the agent's potential to cause harm or engage in undesirable behavior, as discussed in \cite{arxiv:2402.11208}. Ethics-oriented metrics, as highlighted in \cite{arxiv:2405.11357}, consider the agent's adherence to ethical principles, such as fairness, transparency, and accountability.

The development of effective evaluation metrics for large language model-based agents is an active area of research, with several challenges and open questions remaining. For instance, \cite{arxiv:2402.12914} highlights the need for metrics that can capture the nuances of human-agent collaboration, while \cite{doi:10.1145/3748302} emphasizes the importance of evaluating agent memory and its impact on performance. Furthermore, \cite{arxiv:2410.06153} demonstrates the potential of modular design spaces in evaluating and optimizing agent performance.

In conclusion, the evaluation of large language model-based agents requires a multifaceted approach that incorporates various metrics and benchmarks. By synthesizing insights from \cite{arxiv:1906.09379}, \cite{arxiv:2308.03688}, and other relevant studies, researchers can develop a more comprehensive understanding of agent performance and its implications for real-world applications. As the field continues to evolve, it is essential to prioritize the development of robust, generalizable evaluation metrics that can capture the complexities of large language model-based agents and their potential impact on society. Future research directions may include the exploration of novel metrics, such as those that incorporate human evaluation and feedback, as well as the development of more sophisticated benchmarks that can effectively assess agent performance in dynamic, real-world environments, as discussed in \cite{doi:10.18653/v1/2023.acl-long.411} and \cite{arxiv:2306.03917}.


\subsection{Evaluation Frameworks for Large Language Model-based Agents}


The development of comprehensive evaluation frameworks for large language model-based agents is crucial for assessing their performance, identifying areas for improvement, and ensuring their safe and effective deployment in real-world applications. Building on the importance of evaluation metrics and benchmarks discussed in the previous section, it is essential to design evaluation frameworks that can effectively capture the nuances of agent performance, particularly in complex, dynamic environments. As highlighted in \cite{arxiv:2308.03688}, the evaluation of large language model-based agents requires a multi-dimensional approach, considering aspects such as task completion, safety, and ethical considerations. Recent studies, including \cite{doi:10.1145/3748302} and \cite{doi:10.1609/aaaiss.v2i1.27688}, have emphasized the importance of memory mechanisms in large language model-based agents, underscoring the need for evaluation frameworks that account for these mechanisms.

To address the challenges of evaluating large language model-based agents, several evaluation frameworks have been proposed. For instance, \cite{arxiv:2404.06411} introduces a modular benchmark framework that allows for the easy extension of benchmarks and metrics, enabling the evaluation of LLM agents across various tasks and environments. Similarly, \cite{arxiv:2412.14470} provides a comprehensive benchmark for evaluating the safety of LLM agents, covering various safety risks and failure modes. These frameworks demonstrate the growing recognition of the need for specialized evaluation approaches for large language model-based agents, and their development is essential for ensuring the safe and effective deployment of these agents in real-world applications.

The integration of human evaluation and automated metrics is also essential for comprehensive evaluation frameworks. As discussed in \cite{doi:10.1007/s11704-024-40231-1}, human evaluation can provide valuable insights into the performance and limitations of large language model-based agents, while automated metrics can offer a more objective and scalable assessment. The development of hybrid evaluation approaches, combining human and automated assessment, is an active area of research, with studies such as \cite{arxiv:2404.01023} exploring the potential of multi-agent systems for evaluating large language models. By leveraging these hybrid approaches, researchers can create more effective and reliable evaluation frameworks for large language model-based agents.

Emerging trends in evaluation frameworks for large language model-based agents include the use of multimodal evaluation metrics, which can account for the diverse range of tasks and environments in which these agents operate. For example, \cite{arxiv:2411.04890} highlights the importance of evaluating GUI agents in complex, real-world environments, requiring the development of specialized evaluation metrics and frameworks. Additionally, the increasing focus on safety and ethical considerations in large language model-based agents has led to the development of evaluation frameworks that prioritize these aspects, such as \cite{arxiv:2412.14470}. As the field continues to evolve, it is essential to prioritize the development of comprehensive evaluation frameworks that can capture the complexity and nuance of large language model-based agents.

In conclusion, the development of comprehensive evaluation frameworks for large language model-based agents is a critical area of research, requiring a multidisciplinary approach that integrates insights from natural language processing, artificial intelligence, and human-computer interaction. By leveraging advances in evaluation metrics, benchmarking, and human-AI collaboration, researchers can create more effective and reliable evaluation frameworks for large language model-based agents, ultimately contributing to their safe and beneficial deployment in real-world applications. As noted in \cite{arxiv:2501.11354}, the development of robust evaluation frameworks is essential for realizing the full potential of large language model-based agents, and future research should prioritize the creation of more comprehensive, flexible, and scalable evaluation approaches that can inform the optimization and improvement of these agents, as will be discussed in the following section.


\subsection{Challenges and Future Directions in Evaluating Large Language Model-based Agents}


The evaluation of large language model-based agents poses significant challenges due to their complexity, adaptability, and the dynamic nature of their applications. As highlighted in \cite{doi:10.1145/3641289}, the assessment of these models requires a multifaceted approach, considering not only their performance on specific tasks but also their ability to generalize, reason, and interact with humans and other agents in diverse environments. \cite{arxiv:2308.03688} and \cite{arxiv:2308.04026} have made notable contributions by providing benchmarks and simulation environments for evaluating the capabilities of large language model-based agents. However, these efforts also underscore the need for more comprehensive and nuanced evaluation frameworks that can capture the full range of agent behaviors and interactions.

One of the primary challenges in evaluating large language model-based agents is the development of robust and generalizable evaluation metrics. Current metrics, such as perplexity and accuracy, provide limited insights into the agents' ability to reason, solve complex problems, or engage in effective human-computer interaction. \cite{arxiv:2404.09135} emphasizes the importance of metrics that can quantify the agents' performance in real-world scenarios, including their ability to adapt to new tasks, learn from feedback, and demonstrate ethical and safe behavior. 

Another critical aspect of evaluating large language model-based agents is the consideration of their safety and ethical implications. \cite{arxiv:2402.11208} and \cite{arxiv:2412.14470} highlight the potential risks associated with these agents, including backdoor attacks, data poisoning, and unintended consequences of their actions. To mitigate these risks, it is essential to develop evaluation frameworks that incorporate safety and ethical considerations, such as fairness, transparency, and accountability.

The future of evaluating large language model-based agents lies in the development of more sophisticated and comprehensive evaluation frameworks that can capture the complexity and nuance of these systems. \cite{arxiv:2402.11443} and \cite{arxiv:2404.06411} represent significant steps in this direction, as they provide modular and extensible frameworks for evaluating agent performance and safety. 

In conclusion, the evaluation of large language model-based agents is a complex and multifaceted challenge that requires a comprehensive and nuanced approach. By developing more sophisticated evaluation frameworks, incorporating safety and ethical considerations, and creating more realistic simulation environments, researchers can provide a more accurate and comprehensive assessment of these agents' capabilities and limitations. As highlighted in \cite{arxiv:2303.18223}, the future of large language model-based agents depends on the development of effective evaluation methodologies that can ensure their safe, efficient, and beneficial deployment in real-world applications.


\section{Applications and Case Studies of Large Language Model-based Agents}



\subsection{Introduction to Applications and Case Studies}


The emergence of large language model-based agents has revolutionized the field of artificial intelligence, enabling machines to interact with humans in a more natural and intuitive way. These agents have been successfully applied in various domains, including natural language processing, decision-making, and problem-solving, demonstrating their potential to improve efficiency, productivity, and decision-making in numerous industries \cite{doi:10.1145/3604237.3626869}. As highlighted in \cite{doi:10.1007/s11704-024-40231-1}, the optimization of these agents is crucial for enhancing their performance and efficiency. 

One of the key aspects of large language model-based agents is their ability to learn from experience and adapt to new environments. This is particularly important in applications such as customer service, where agents need to respond to a wide range of queries and provide personalized support \cite{doi:10.1145/3624724}. The use of large language models in these applications has been shown to improve customer satisfaction and reduce response times \cite{doi:10.1109/mnet.2024.3435752}. Furthermore, the integration of large language models with other technologies, such as computer vision and robotics, has the potential to enable more sophisticated and human-like interactions \cite{arxiv:2401.13601}.

Despite the many successes of large language model-based agents, there are still several challenges that need to be addressed. One of the main limitations of these agents is their lack of common sense and real-world experience, which can lead to unrealistic or inappropriate responses \cite{doi:10.1016/j.tics.2023.08.006}. Additionally, the training data used to develop these agents can be biased or incomplete, which can result in unfair or discriminatory outcomes \cite{doi:10.1007/s10849-023-09409-x}. To overcome these challenges, it is essential to develop more robust and transparent evaluation metrics, as well as more effective methods for training and fine-tuning large language models \cite{doi:10.1145/3641289}.

The applications of large language model-based agents are diverse and continue to expand. In the field of finance, for example, these agents can be used to analyze financial data, predict market trends, and provide personalized investment advice \cite{doi:10.1145/3604237.3626869}. In healthcare, they can be used to analyze medical records, diagnose diseases, and develop personalized treatment plans \cite{doi:10.1007/s13042-024-02318-w}. In education, they can be used to develop personalized learning plans, provide real-time feedback, and enhance student engagement \cite{doi:10.1145/3649506}.

In conclusion, large language model-based agents have the potential to revolutionize numerous industries and improve our daily lives. However, to realize this potential, it is essential to address the challenges associated with these agents, including their lack of common sense and real-world experience, as well as the biases and limitations of their training data. By developing more robust and transparent evaluation metrics, and more effective methods for training and fine-tuning large language models, we can unlock the full potential of these agents and create more sophisticated and human-like interactions \cite{doi:10.3390/app14052074}. As noted in \cite{arxiv:2206.07682}, the continued development and application of large language model-based agents will depend on our ability to understand and address these challenges, and to develop more effective methods for evaluating and improving their performance. 

The future of large language model-based agents is exciting and uncertain, with many potential applications and challenges on the horizon. As highlighted in \cite{doi:10.34133/research.0655}, the development of these agents will depend on our ability to balance their potential benefits with their potential risks, and to ensure that they are developed and used in a responsible and transparent way. By working together to address these challenges, we can unlock the full potential of large language model-based agents and create a brighter future for all \cite{arxiv:2404.01023}. Ultimately, the success of these agents will depend on our ability to develop more sophisticated and human-like interactions, and to create more robust and transparent evaluation metrics \cite{arxiv:2402.10172}. As noted in \cite{arxiv:2406.04151}, the development of these agents is an ongoing process, and their potential applications and challenges will continue to evolve in the coming years.


\subsection{Natural Language Processing Applications}


The applications of large language model-based agents in natural language processing (NLP) have been rapidly expanding, with significant advancements in tasks such as text classification, sentiment analysis, and machine translation. As discussed in the previous section, large language model-based agents have the potential to revolutionize numerous industries and improve our daily lives, and their application in NLP tasks is a crucial aspect of this potential. \cite{arxiv:2309.14365} highlights the importance of optimization techniques in enhancing the performance of large language models, which is crucial for their application in NLP tasks.

One of the key areas where large language model-based agents have shown impressive performance is in text classification. \cite{doi:10.1145/3605943} provides an overview of the current state of pre-trained language models and their applications in NLP tasks, including text classification. The use of large language models for text classification has been shown to achieve state-of-the-art results, outperforming traditional machine learning approaches. \cite{arxiv:1412.6980} presents a stochastic optimization method that can be used to fine-tune large language models for text classification tasks. Furthermore, the integration of large language models with other technologies, such as computer vision and robotics, has the potential to enable more sophisticated and human-like interactions, as discussed in \cite{arxiv:2401.13601}.

Sentiment analysis is another NLP task where large language model-based agents have demonstrated significant capabilities. \cite{arxiv:2402.06196} discusses the applications of large language models in sentiment analysis, including their ability to capture nuanced aspects of language and provide more accurate predictions. Additionally, the use of large language models in sentiment analysis has been shown to improve customer satisfaction and reduce response times, as highlighted in \cite{doi:10.1109/mnet.2024.3435752}.

Machine translation is another area where large language model-based agents have shown remarkable performance. \cite{arxiv:2402.16142} provides an overview of the applications of large language models in machine translation, including their ability to capture complex linguistic structures and provide more accurate translations. \cite{arxiv:2403.19887} presents a hybrid language model that combines the strengths of transformer and Mamba models for machine translation tasks. The development of more robust and transparent evaluation metrics, as well as more effective methods for training and fine-tuning large language models, is crucial for improving their performance in machine translation tasks, as discussed in \cite{doi:10.1145/3641289}.

Despite the significant advancements in NLP tasks, there are still several challenges that need to be addressed. One of the key challenges is the lack of interpretability and explainability of large language models. \cite{doi:10.1007/s10849-023-09409-x} discusses the strengths and weaknesses of large language models, including their lack of interpretability and explainability. Addressing these challenges will be crucial to fully realizing the potential of large language model-based agents in NLP tasks and decision-making applications, as discussed in the following section.

In conclusion, large language model-based agents have demonstrated significant capabilities in NLP tasks, including text classification, sentiment analysis, and machine translation. However, there are still several challenges that need to be addressed, including the lack of interpretability and explainability, and the need for large amounts of training data. \cite{doi:10.3390/app14052074} discusses the current trends and challenges in large language models. \cite{arxiv:2402.14744} presents a novel approach to using large language models for personal mobility generation. Overall, the application of large language model-based agents in NLP tasks has the potential to revolutionize the field and improve decision-making capabilities in various domains, leading to new opportunities for growth and innovation in the field of artificial intelligence.


\subsection{Decision-Making and Problem-Solving Applications}


The application of large language model-based agents in decision-making and problem-solving tasks has garnered significant attention in recent years. These agents, equipped with advanced natural language processing capabilities, can process and generate human-like language, enabling them to interact with humans and other agents in complex environments \cite{arxiv:1906.09379}. The use of large language models in decision-making and problem-solving tasks is particularly promising, as they can leverage their language understanding capabilities to analyze complex situations, identify patterns, and make informed decisions \cite{doi:10.18653/v1/2023.acl-long.411}.

One of the primary applications of large language model-based agents in decision-making is in expert systems. These systems, designed to mimic the decision-making abilities of human experts, can be used in a variety of domains, including healthcare, finance, and education \cite{arxiv:2304.05332}. For instance, a large language model-based agent can be used to analyze medical records, identify potential diagnoses, and recommend treatment options \cite{doi:10.1038/s41746-024-01083-y}. Similarly, in finance, these agents can be used to analyze market trends, identify potential investment opportunities, and make recommendations to investors \cite{arxiv:2408.06361}.

In addition to expert systems, large language model-based agents can also be used in planning and optimization tasks. These tasks, which involve identifying the most effective course of action to achieve a specific goal, can be particularly challenging in complex environments \cite{arxiv:2205.05718}. Large language model-based agents can be used to analyze the environment, identify potential obstacles, and develop plans to overcome them \cite{arxiv:2305.17144}. For example, a large language model-based agent can be used to plan a route for a self-driving car, taking into account traffic patterns, road conditions, and other factors \cite{doi:10.1109/mnet.2024.3435752}.

Despite the potential of large language model-based agents in decision-making and problem-solving tasks, there are several challenges that must be addressed. One of the primary challenges is the need for these agents to be able to reason and make decisions in a way that is transparent and explainable \cite{arxiv:2306.03917}. This requires the development of new techniques for interpreting and understanding the decisions made by these agents \cite{arxiv:2308.03688}. Another challenge is the need for these agents to be able to learn from feedback and adapt to new environments \cite{arxiv:2308.05960}.

In recent years, there have been several studies that have explored the use of large language model-based agents in decision-making and problem-solving tasks. For example, \cite{arxiv:2308.04026} presented a sandbox environment for evaluating the decision-making capabilities of large language model-based agents. Similarly, \cite{arxiv:2402.01680} presented a survey of the use of large language model-based agents in multi-agent systems. These studies have demonstrated the potential of large language model-based agents in decision-making and problem-solving tasks, but have also highlighted the need for further research in this area \cite{arxiv:2401.02777}.

In conclusion, the application of large language model-based agents in decision-making and problem-solving tasks has the potential to revolutionize the way we approach complex problems. These agents, equipped with advanced natural language processing capabilities, can process and generate human-like language, enabling them to interact with humans and other agents in complex environments. However, there are several challenges that must be addressed, including the need for transparency and explainability in decision-making, and the need for these agents to be able to learn from feedback and adapt to new environments. Further research is needed to fully realize the potential of large language model-based agents in decision-making and problem-solving tasks \cite{arxiv:2406.04692}. As noted in \cite{arxiv:2404.01023}, the use of multiple agents can enhance the decision-making capabilities of large language models, and \cite{arxiv:2407.07086} presented a framework for scaffolding theory of mind in multi-agent tasks. Overall, the use of large language model-based agents in decision-making and problem-solving tasks is a rapidly evolving field, with significant potential for innovation and advancement \cite{arxiv:2303.18223}.


\subsection{Human-Computer Interaction and Collaboration}


The application of large language model-based agents in human-computer interaction and collaboration has been a rapidly evolving field, with significant advancements in recent years, building on the potential of these agents in decision-making and problem-solving tasks discussed earlier. As noted in \cite{arxiv:2206.07682}, the capabilities of large language models have improved substantially, enabling them to perform complex tasks and interact with humans in a more natural and intuitive way. One of the primary areas where large language model-based agents have shown great promise is in chatbots and virtual assistants, where they can understand and respond to user queries, provide information, and even perform tasks on behalf of the user \cite{arxiv:2309.07870}. 

The use of large language model-based agents in collaborative problem-solving has also been explored, where multiple agents can work together to achieve a common goal \cite{doi:10.1145/3712003}. This approach has shown significant potential in improving the efficiency and effectiveness of problem-solving, as well as enhancing the overall user experience. For instance, \cite{arxiv:2308.03688} presents a multi-dimensional benchmark for evaluating the performance of large language model-based agents, highlighting the importance of considering multiple facets of agent performance, including efficiency, effectiveness, and safety. Furthermore, the application of large language model-based agents in human-computer interaction and collaboration is closely related to their evaluation and benchmarking, which is crucial for assessing their performance and identifying areas for improvement, as will be discussed in the next section.

In addition to their application in chatbots and collaborative problem-solving, large language model-based agents have also been used in various other areas, such as customer service, tech support, and language translation \cite{arxiv:2411.04890}. These agents have been shown to be highly effective in providing personalized support and assistance to users, and have the potential to revolutionize the way we interact with computers and other devices. As \cite{arxiv:2303.18223} notes, the development of large language models has been driven by the need for more advanced and sophisticated natural language processing capabilities, and has led to significant advancements in areas such as language understanding, generation, and translation.

However, despite the many advantages of large language model-based agents, there are also several challenges and limitations that need to be addressed. One of the primary challenges is ensuring the safety and security of these agents, as they can potentially be used to spread misinformation or engage in malicious activities \cite{arxiv:2402.11208}. Another challenge is ensuring that these agents are transparent and explainable, so that users can understand how they are making decisions and taking actions \cite{arxiv:2402.01108}. 

To address these challenges, researchers have proposed various approaches, including the use of formal methods and techniques to ensure the safety and security of large language model-based agents \cite{arxiv:2402.00798}. Others have proposed the use of multi-agent systems, where multiple agents can work together to achieve a common goal, and can provide a more robust and resilient approach to problem-solving \cite{arxiv:2402.01680}. As \cite{arxiv:2410.06153} notes, the development of large language model-based agents requires a modular and flexible approach, where agents can be easily integrated and composed to achieve complex tasks.

In conclusion, the application of large language model-based agents in human-computer interaction and collaboration has shown significant promise, with potential applications in areas such as chatbots, virtual assistants, and collaborative problem-solving. However, there are also several challenges and limitations that need to be addressed, including ensuring the safety and security of these agents, and ensuring that they are transparent and explainable. Further research is needed to address these challenges, and to fully realize the potential of large language model-based agents in human-computer interaction and collaboration, ultimately leading to more efficient and effective human-computer interaction and collaboration systems, as \cite{arxiv:2501.11354} notes, the development of large language model-based agents requires a long-term vision and roadmap, where researchers and practitioners can work together to address the challenges and limitations of these agents.


\subsection{Evaluation and Benchmarking of Large Language Model-Based Agents}


The evaluation and benchmarking of large language model-based agents are crucial steps in assessing their performance, identifying areas for improvement, and ensuring their safe and effective deployment in real-world applications. As highlighted in \cite{arxiv:2308.03688}, the development of comprehensive evaluation metrics and benchmarks is essential for comparing the performance of different large language model-based agents and identifying areas for improvement. \cite{doi:10.1145/3748302} emphasizes the importance of evaluating the memory mechanisms of large language model-based agents, which play a critical role in their ability to process and generate human-like language.

Recent studies, such as \cite{arxiv:2410.06153} and \cite{arxiv:2406.04151}, have proposed novel approaches to evaluating and benchmarking large language model-based agents. These approaches include the use of modular design spaces, evolutionary algorithms, and multi-agent systems to assess the performance of large language model-based agents in diverse environments and tasks. \cite{arxiv:2402.11443} introduces a multi-agent framework for dynamically evaluating large language models, which enables the assessment of their capabilities and limitations in a more accurate and comprehensive manner.

The evaluation and benchmarking of large language model-based agents also involve assessing their safety and robustness. \cite{arxiv:2412.14470} proposes a comprehensive benchmark for evaluating the safety of large language model-based agents, which includes 349 interaction environments and 2,000 test cases. \cite{arxiv:2402.11208} highlights the importance of investigating backdoor threats to large language model-based agents, which can compromise their safety and security.

In addition to these efforts, researchers have also explored the use of large language models as optimizers for hyperparameter tuning and other optimization tasks. \cite{arxiv:2402.01881} proposes a novel paradigm for hyperparameter optimization using large language models, which has shown promising results in improving the performance of machine learning models. \cite{doi:10.1109/access.2024.3524176} introduces a hybrid approach that combines large language models with metaheuristics for optimization, which has demonstrated significant improvements in solution quality and efficiency.

The development of comprehensive evaluation metrics and benchmarks for large language model-based agents is an ongoing research area. \cite{doi:10.1145/3641289} provides a comprehensive review of existing evaluation methods and metrics for large language models, highlighting their strengths and limitations. \cite{arxiv:2404.09135} emphasizes the importance of evaluating large language models using a range of metrics, including those that assess their safety, robustness, and fairness.

In conclusion, the evaluation and benchmarking of large language model-based agents are critical components of their development and deployment. Recent studies have proposed novel approaches to evaluating and benchmarking these agents, including the use of modular design spaces, evolutionary algorithms, and multi-agent systems. The development of comprehensive evaluation metrics and benchmarks is an ongoing research area, with a focus on assessing the safety, robustness, and fairness of large language model-based agents. As highlighted in \cite{arxiv:2405.19616}, the evaluation and benchmarking of large language model-based agents are essential for identifying their limitations and areas for improvement, which is critical for ensuring their safe and effective deployment in real-world applications. Future research should continue to explore novel approaches to evaluating and benchmarking large language model-based agents, with a focus on developing more comprehensive and robust evaluation metrics and benchmarks.


\subsection{Future Directions and Challenges}


The development and application of large language model-based agents have witnessed significant advancements in recent years, with these agents demonstrating remarkable capabilities in natural language understanding, generation, and decision-making. As a natural progression from the evaluation and benchmarking of large language model-based agents, it is essential to address the challenges and future directions that remain to be explored. One of the primary concerns is the need for continued research in optimization techniques, as highlighted in \cite{arxiv:2312.00678}, to improve the efficiency and scalability of large language models. This is crucial for deploying these models in real-world applications, where computational resources and latency are critical factors.

Building on the importance of evaluation frameworks, as discussed in the previous section, the development of comprehensive evaluation frameworks and benchmarks is essential for assessing the performance of large language model-based agents in various tasks and environments. Furthermore, the integration of ethical considerations, such as fairness, transparency, and accountability, is vital for ensuring the responsible development and deployment of large language model-based agents, as emphasized in \cite{arxiv:2402.06196} and \cite{doi:10.1109/access.2024.3365742}. The evaluation of large language model-based agents, as discussed in \cite{arxiv:2308.03688} and \cite{doi:10.1038/s41746-024-01083-y}, is a critical step in identifying areas for improvement and ensuring their safe and effective deployment.

In addition to these challenges, the application of large language model-based agents in emerging areas, such as edge computing and human-computer interaction, presents significant opportunities for future research. For instance, \cite{doi:10.1109/mnet.2025.3539338} explores the potential of large language models in edge computing, while \cite{doi:10.1109/mwc.005.2400019} discusses their application in 6G networks. The use of large language models in multi-agent systems, as presented in \cite{arxiv:2402.01680} and \cite{doi:10.1145/3712003}, also offers a promising direction for future research. These emerging areas and applications will be crucial in shaping the future of large language model-based agents and their potential impact on various aspects of society.

Moreover, the development of more advanced architectures, such as multimodal large language models, is expected to enhance the capabilities of large language model-based agents. As discussed in \cite{arxiv:2401.13601}, multimodal large language models can process and generate multiple forms of data, including text, images, and audio, enabling more versatile and human-like interactions. The application of large language models in specific domains, such as healthcare and finance, also presents significant opportunities for future research, as highlighted in \cite{arxiv:2408.06361} and \cite{doi:10.3348/kjr.2023.0997}. These advancements and applications will be essential in driving the growth and adoption of large language model-based agents in various industries and domains.

However, despite these opportunities, several challenges need to be addressed, including the lack of interpretability, the risk of hallucinations, and the need for more efficient training methods. As discussed in \cite{arxiv:2405.03207}, the development of more interpretable and explainable large language models is crucial for understanding their decision-making processes and ensuring their reliability. Furthermore, the application of large language models in real-world scenarios requires careful consideration of their limitations and potential biases, as emphasized in \cite{arxiv:2402.14744} and \cite{arxiv:2401.02777}. By acknowledging and addressing these challenges, researchers and practitioners can work together to create more advanced, efficient, and responsible large language model-based agents.

In conclusion, the development and application of large language model-based agents present significant opportunities for future research, with potential applications in emerging areas, such as edge computing, human-computer interaction, and multi-agent systems. However, addressing the challenges associated with optimization, evaluation, and ethical considerations is crucial for ensuring the responsible development and deployment of these agents. As the field continues to evolve, it is essential to prioritize the development of more advanced, efficient, and responsible large language model-based agents that can positively impact various aspects of society, which will be further explored in the following sections.


\section{Safety, Ethics, and Regulatory Considerations}



\subsection{Introduction to Safety, Ethics, and Regulatory Considerations}


The development and deployment of large language model-based agents have transformed the landscape of artificial intelligence, with profound implications for various aspects of society. As these agents become increasingly integrated into our daily lives, it is essential to consider the safety, ethics, and regulatory considerations that arise from their use \cite{doi:10.1109/tkde.2025.3554028}. The importance of responsible AI practices cannot be overstated, as the potential risks and consequences of unethical or unregulated large language model-based agents are significant \cite{arxiv:2402.06196}. 

One of the primary concerns is the potential for bias and discrimination in large language models, which can perpetuate existing social inequalities \cite{doi:10.1007/s10849-023-09409-x}. For instance, a study by \cite{arxiv:2212.09251} found that large language models can exhibit biased behavior, such as generating text that reflects existing social biases. To mitigate these risks, it is crucial to develop and implement robust evaluation metrics and frameworks that can detect and address bias in large language models \cite{doi:10.1145/3641289}. 

Another critical aspect is the need for transparency and explainability in large language models \cite{doi:10.1016/j.tics.2023.08.006}. As these models become increasingly complex, it is essential to develop techniques that can provide insights into their decision-making processes and ensure that their outputs are fair, reliable, and trustworthy \cite{arxiv:2311.07226}. Furthermore, the integration of large language models with other technologies, such as computer vision and robotics, raises additional safety and ethical concerns that must be addressed \cite{doi:10.1145/3604237.3626869}. 

Regulatory frameworks play a vital role in ensuring the safe and responsible development and deployment of large language model-based agents \cite{arxiv:2309.14365}. Governments and regulatory bodies must establish clear guidelines and standards for the development and use of these agents, including requirements for transparency, accountability, and fairness \cite{doi:10.3390/app14052074}. Additionally, there is a need for ongoing research and development in the field of large language models, including the exploration of new architectures, training methods, and evaluation metrics \cite{arxiv:2402.01364}. 

The potential applications of large language model-based agents are vast and varied, ranging from natural language processing and decision-making to problem-solving and human-computer interaction \cite{arxiv:2309.02427}. However, as these agents become increasingly autonomous and integrated into complex systems, there is a growing need for robust safety and ethical protocols to prevent potential risks and harms \cite{doi:10.1038/s41746-024-01083-y}. The development of large language model-based agents that can learn from feedback, adapt to new environments, and make decisions in real-time raises significant challenges and opportunities for the field of artificial intelligence \cite{arxiv:2402.01680}. 

In conclusion, the development and deployment of large language model-based agents require a comprehensive and multidisciplinary approach that prioritizes safety, ethics, and regulatory considerations \cite{doi:10.1109/mnet.2024.3435752}. By acknowledging the potential risks and challenges associated with these agents and working to develop robust evaluation metrics, frameworks, and regulatory protocols, we can ensure that large language model-based agents are developed and used in a responsible and beneficial manner \cite{arxiv:2405.03207}. As the field of artificial intelligence continues to evolve, it is essential to prioritize transparency, accountability, and fairness in the development and deployment of large language model-based agents, and to recognize the critical role that these agents will play in shaping the future of our society \cite{doi:10.54254/2755-2721/35/20230399}.


\subsection{Ethical Considerations for Large Language Model-based Agents}


The development and deployment of large language model-based agents have raised significant ethical considerations, including ensuring transparency, fairness, and respect for user autonomy \cite{arxiv:2309.14365}. As these agents become increasingly integrated into various aspects of life, from customer service to healthcare, it is crucial to address the potential risks and consequences of their actions \cite{arxiv:2402.06196}. One of the primary concerns is the lack of transparency in decision-making processes, which can lead to unintended biases and discrimination \cite{doi:10.1007/s10849-023-09409-x}. For instance, a language model may perpetuate existing social biases if it is trained on biased data, resulting in unfair outcomes \cite{arxiv:2311.07621}. This concern is closely tied to the need for responsible AI practices, as highlighted earlier, and underscores the importance of developing and implementing robust evaluation metrics and frameworks that can detect and address bias in large language models \cite{doi:10.1145/3641289}.

To mitigate these risks, researchers have proposed various approaches, including data curation and debiasing techniques \cite{doi:10.1145/3605943}. Additionally, the development of explainable AI methods can help provide insights into the decision-making processes of large language models, enabling more transparent and accountable systems \cite{doi:10.3390/app14052074}. However, these methods are not without their limitations, and further research is needed to address the complex ethical challenges associated with large language model-based agents \cite{arxiv:2206.07682}. The importance of transparency and explainability in large language models is also closely tied to the need for regulatory frameworks that prioritize safety, ethics, and accountability, as discussed earlier \cite{arxiv:2309.14365}.

Another critical aspect is ensuring fairness and respect for user autonomy. Large language models can be used to manipulate or deceive users, which can have severe consequences, particularly in high-stakes applications such as finance or healthcare \cite{arxiv:2408.06361}. To prevent such misuse, it is essential to develop and implement robust evaluation metrics and benchmarks that assess the performance of large language models in terms of fairness, transparency, and accountability \cite{arxiv:2308.03688}. Furthermore, researchers have proposed the use of value alignment techniques, which aim to align the objectives of large language models with human values, promoting more ethical and responsible decision-making \cite{arxiv:2405.03207}. This is particularly important as large language models become increasingly integrated into complex systems, where their decisions can have significant consequences \cite{arxiv:2311.07226}.

The integration of large language models with other technologies, such as computer vision or robotics, also raises important ethical considerations \cite{arxiv:2402.16142}. For example, the use of large language models in autonomous vehicles or drones can have significant consequences, particularly if these systems are not designed with safety and ethics in mind \cite{arxiv:2402.14744}. Therefore, it is crucial to develop and implement robust safety protocols and ethical guidelines that ensure the responsible development and deployment of large language model-based agents \cite{doi:10.34133/research.0655}. This will be further explored in the context of safety measures and regulatory frameworks for large language model-based agents, highlighting the need for a comprehensive and multidisciplinary approach to addressing these challenges \cite{arxiv:2402.06196}.

In conclusion, the ethical considerations associated with large language model-based agents are complex and multifaceted, requiring careful attention to transparency, fairness, and respect for user autonomy \cite{arxiv:2303.18223}. As researchers continue to develop and deploy these systems, it is essential to prioritize ethical considerations and develop robust evaluation metrics and benchmarks that assess their performance in terms of fairness, transparency, and accountability \cite{arxiv:2410.02644}. By doing so, we can ensure that large language model-based agents are developed and deployed in a responsible and ethical manner, promoting beneficial outcomes for individuals and society as a whole \cite{arxiv:2404.01023}. Ultimately, this will require ongoing research and development in the field of large language models, as well as collaboration between researchers, policymakers, and industry stakeholders to establish clear guidelines and standards for the development and use of these agents \cite{arxiv:2402.01364}.


\subsection{Safety Measures for Large Language Model-based Agents}


The development and deployment of large language model-based agents have raised significant concerns regarding their safety and potential misuse. As these agents become increasingly sophisticated and autonomous, it is essential to implement robust safety measures to prevent harmful behaviors \cite{doi:10.1038/s41746-024-01083-y}. One critical approach is to design mechanisms for detecting and mitigating harmful actions, such as those that could cause physical harm or perpetuate misinformation \cite{arxiv:2402.11208}. 

Researchers have proposed various methods to address these concerns, including the development of formal frameworks for evaluating and benchmarking the safety of large language model-based agents \cite{arxiv:2410.02644}. These frameworks can help identify potential vulnerabilities and provide a foundation for developing more secure and reliable agents. Furthermore, techniques such as reinforcement learning from human feedback \cite{doi:10.18653/v1/2021.naacl-main.341} and multi-agent collaboration \cite{arxiv:2402.12914} can be used to improve the safety and performance of large language model-based agents.

Another crucial aspect of ensuring safety is to develop agents that can reason about their own limitations and uncertainties \cite{arxiv:2205.05718}. This can be achieved through the integration of cognitive architectures and large language models, enabling agents to better understand their own capabilities and potential biases \cite{doi:10.1609/aaaiss.v2i1.27691}. Additionally, the use of external knowledge sources and memory mechanisms can help agents to better contextualize their actions and decisions \cite{doi:10.1609/aaaiss.v2i1.27688}.

The importance of evaluating and assessing the safety of large language model-based agents cannot be overstated \cite{arxiv:2308.04026}. This requires the development of comprehensive evaluation metrics and benchmarks that can effectively measure the performance and safety of these agents in various scenarios \cite{arxiv:2308.05960}. Moreover, the need for transparency and explainability in the decision-making processes of large language model-based agents is critical for building trust and ensuring safety \cite{arxiv:2306.03917}.

In conclusion, the safety of large language model-based agents is a complex and multifaceted issue that requires careful consideration and addressing. By developing and implementing robust safety measures, such as mechanisms for detecting and mitigating harmful behaviors, and evaluating and assessing the safety of these agents, we can work towards creating more reliable and trustworthy autonomous systems \cite{doi:10.1109/mwc.016.2300600}. As research in this area continues to evolve, it is essential to prioritize the development of safe and responsible large language model-based agents that can benefit society while minimizing potential risks \cite{doi:10.18653/v1/2023.findings-emnlp.461}. Future directions in this field should focus on addressing the challenges and limitations of current safety measures, such as the need for more effective evaluation metrics and benchmarks, and exploring new approaches to improving the safety and reliability of large language model-based agents \cite{arxiv:2410.06153}.


\subsection{Regulatory Frameworks for Large Language Model-based Agents}


The development and deployment of large language model-based agents have sparked intense interest in recent years, with their potential applications spanning various domains, including natural language processing, decision-making, and problem-solving. As these agents become increasingly integrated into our daily lives, it is crucial to establish regulatory frameworks that ensure their safe and responsible development and deployment. Building on the safety concerns discussed in the previous section, this subsection delves into the regulatory frameworks governing large language model-based agents, including guidelines for compliance and enforcement, as discussed in \cite{doi:10.1145/3748302} and \cite{doi:10.1145/3712003}.

One of the primary challenges in regulating large language model-based agents is the lack of standardization in their development and deployment. As noted in \cite{arxiv:2308.03688}, the absence of a unified framework for evaluating and assessing these agents hinders the establishment of effective regulatory measures. To address this, researchers have proposed various benchmarks and evaluation metrics, such as those presented in \cite{arxiv:2404.06411} and \cite{arxiv:2412.14470}. These benchmarks and metrics provide a foundation for assessing the performance and safety of large language model-based agents, enabling regulators to develop more effective guidelines for their development and deployment. Furthermore, the development of comprehensive evaluation frameworks, as discussed in the following section, is essential for ensuring the safety and regulatory compliance of large language model-based agents.

Another critical aspect of regulating large language model-based agents is ensuring their compliance with existing laws and regulations. As discussed in \cite{doi:10.18653/v1/2023.findings-emnlp.461}, large language model-based agents must comply with regulations related to data privacy, security, and intellectual property. To achieve this, developers and regulators must work together to establish clear guidelines and standards for the development and deployment of these agents. For instance, \cite{arxiv:2308.04030} proposes a collaborative platform for developing and sharing large language model-based agents, which can facilitate the establishment of common standards and guidelines for their development and deployment.

In addition to ensuring compliance, the regulatory frameworks governing large language model-based agents must also address the issues of accountability and explainability. As noted in \cite{arxiv:2402.11208}, large language model-based agents can be vulnerable to backdoor attacks, which can compromise their safety and security. To mitigate this risk, regulators must establish guidelines for ensuring the accountability of developers and deployers of large language model-based agents. This can include measures such as transparency in development and deployment practices, regular auditing and testing, and the establishment of clear lines of responsibility in case of errors or incidents. Moreover, as discussed in \cite{arxiv:2402.01108}, large language model-based agents can be difficult to interpret and understand, which can hinder their adoption in critical applications. To address this, regulators must establish guidelines for ensuring the explainability of large language model-based agents, such as the use of transparent and interpretable models, and the provision of clear explanations for their decisions and actions.

In conclusion, the regulatory frameworks governing large language model-based agents are complex and multifaceted, requiring careful consideration of various factors, including standardization, compliance, accountability, and explainability. As noted in \cite{arxiv:2303.18223}, the development and deployment of large language model-based agents have the potential to transform various domains, but their safe and responsible development and deployment require effective regulatory measures. By establishing clear guidelines and standards for the development and deployment of large language model-based agents, regulators can ensure their safe and beneficial use, while also promoting innovation and progress in the field. The evaluation and assessment of large language model-based agents, as discussed in the following section, play a crucial role in ensuring their safety, ethics, and regulatory compliance, and should be considered in conjunction with the regulatory frameworks governing these agents.


\subsection{Evaluation and Assessment of Large Language Model-based Agents}


The evaluation and assessment of large language model-based agents are crucial for ensuring their safety, ethics, and regulatory compliance. As these agents become increasingly integrated into various aspects of our lives, it is essential to develop comprehensive benchmarks and standards to measure their performance, reliability, and potential risks \cite{arxiv:2308.03688}. The development of such benchmarks is a complex task, requiring a multidisciplinary approach that involves expertise in natural language processing, cognitive science, ethics, and regulatory affairs \cite{doi:10.1145/3748302}.

Recent studies have proposed various evaluation frameworks for large language model-based agents, including AgentBench \cite{arxiv:2308.03688} and AgentQuest \cite{arxiv:2404.06411}. These frameworks provide a comprehensive set of tasks and metrics to assess the agents' performance, safety, and ethics. For instance, AgentBench evaluates agents' ability to complete tasks, follow instructions, and avoid harmful behaviors, while AgentQuest focuses on the agents' ability to reason, learn, and adapt to new situations \cite{arxiv:2404.06411}.

However, the development of benchmarks for large language model-based agents is still in its early stages, and there are several challenges that need to be addressed. One of the primary challenges is the lack of standardization in the evaluation of these agents, which makes it difficult to compare their performance and identify areas for improvement \cite{doi:10.1145/3641289}. Another challenge is the need for more nuanced and context-dependent evaluation frameworks that can capture the complexities of human-language interaction and the potential risks associated with these agents \cite{arxiv:2404.09135}.

To address these challenges, researchers have proposed various approaches, including the use of multimodal evaluation frameworks \cite{arxiv:2406.04692} and the development of more advanced metrics that can capture the agents' performance, safety, and ethics \cite{arxiv:2405.19616}. For example, the use of multimodal evaluation frameworks can help to assess the agents' ability to understand and generate human-like language, as well as their ability to reason and learn in complex environments \cite{arxiv:2406.04692}.

In addition to the development of benchmarks and evaluation frameworks, it is essential to consider the ethical and regulatory implications of large language model-based agents. These agents have the potential to impact various aspects of our lives, from education and healthcare to finance and transportation, and it is crucial to ensure that they are designed and developed with ethical and regulatory considerations in mind \cite{arxiv:2402.11208}. This includes ensuring that the agents are transparent, fair, and accountable, and that they do not perpetuate biases or discriminate against certain groups of people \cite{doi:10.1007/s11704-024-40231-1}.

In conclusion, the evaluation and assessment of large language model-based agents are critical for ensuring their safety, ethics, and regulatory compliance. The development of comprehensive benchmarks and standards is essential for measuring their performance, reliability, and potential risks. While there are several challenges that need to be addressed, researchers are making progress in developing more nuanced and context-dependent evaluation frameworks, as well as considering the ethical and regulatory implications of these agents \cite{doi:10.1145/3748302}. As the field continues to evolve, it is essential to prioritize the development of benchmarks and evaluation frameworks that can capture the complexities of human-language interaction and the potential risks associated with large language model-based agents \cite{arxiv:2412.14470}.


\subsection{Future Directions and Challenges}


As the field of large language model-based agents continues to evolve, several future directions and challenges emerge that warrant attention from researchers and practitioners alike, building on the evaluation frameworks and regulatory considerations discussed earlier \cite{arxiv:2402.06196}. One of the primary concerns is the need for ongoing research and development in ensuring the safety, ethics, and regulatory compliance of these agents, which is crucial for their widespread adoption and deployment \cite{doi:10.3390/app14052074}. For instance, the development of more robust evaluation frameworks that can assess the performance of large language model-based agents in dynamic and real-world environments is essential, as highlighted by the limitations of current evaluation frameworks \cite{doi:10.1038/s41746-024-01083-y}.

Another significant challenge is the integration of large language models with other technologies, such as computer vision and edge computing, to create more comprehensive and interactive agent systems, which requires advances in areas like multimodal learning and human-computer interaction \cite{doi:10.1109/mwc.005.2400019}. This integration also raises important questions about data privacy, security, and regulatory compliance, particularly in domains like finance and healthcare where large language model-based agents are being increasingly applied \cite{arxiv:2408.06361}. Furthermore, the development of new architectures and algorithms that can efficiently process and generate multimodal data is necessary to support the creation of more advanced agent systems \cite{arxiv:2401.13601}.

The use of large language models as agents also poses challenges related to transparency, explainability, and accountability, which are critical for ensuring that these agents are aligned with human values and ethics \cite{doi:10.1145/3748302}. As these models become more complex and autonomous, it is essential to develop methods that can provide insights into their decision-making processes and ensure that they are fair and transparent \cite{doi:10.18653/v1/2023.acl-long.411}. This may involve the development of new techniques for interpretability and explainability, as well as the creation of frameworks that can evaluate the fairness and transparency of large language model-based agents, such as AgentBench \cite{arxiv:2308.03688}.

In addition to these technical challenges, there are also important social and ethical considerations that must be addressed, particularly in areas like education and employment where large language model-based agents are being increasingly deployed \cite{doi:10.1145/3712003}. It is essential to develop guidelines and regulations that can ensure the responsible development and deployment of these agents, as well as provide protections for individuals and communities that may be impacted by their use \cite{doi:10.1109/mnet.2025.3539338}. By prioritizing these considerations, researchers and practitioners can work towards creating a future where large language model-based agents can be used to benefit society as a whole, which will be discussed in more detail in the following section \cite{doi:10.34133/research.0655}.

Finally, the future of large language model-based agents will depend on the ability of researchers and practitioners to address these challenges and develop innovative solutions that can harness the potential of these technologies while minimizing their risks \cite{arxiv:2408.15769}. This will require ongoing investment in research and development, as well as collaboration between academia, industry, and government to ensure that the benefits of these technologies are equitably distributed and that their risks are mitigated \cite{arxiv:2402.12914}. By working together to address these challenges, we can unlock the full potential of large language model-based agents and create a future where these technologies can be used to drive positive change and improvement in various aspects of our lives.


\section{Conclusion}


The development and application of large language model-based agents have witnessed significant advancements in recent years, with these models demonstrating unprecedented capabilities in natural language understanding and generation. As highlighted in \cite{arxiv:2206.07682}, the scaling up of language models has led to the emergence of novel abilities that are not present in smaller models, underscoring the importance of continued research into the optimization and evaluation of these agents. The survey of \cite{arxiv:2309.14365} provides a comprehensive overview of the current state of the field, emphasizing the need for innovative approaches to optimizing large language models and highlighting the potential applications of these models in various domains.

The integration of large language models with other techniques, such as reinforcement learning and multi-task learning, has been shown to enhance their performance and versatility. For instance, \cite{doi:10.18653/v1/2023.acl-long.411} discusses the application of large language models in code generation and programming tasks, demonstrating their potential in automating software development and improving coding efficiency. Similarly, \cite{arxiv:2402.01680} explores the use of large language models in multi-agent systems, highlighting their ability to facilitate complex decision-making and cooperation in dynamic environments.

Despite the significant progress made in the development of large language model-based agents, several challenges and limitations remain to be addressed. As noted in \cite{doi:10.1007/s10849-023-09409-x}, these models can be prone to bias, hallucinations, and lack of transparency, which can limit their reliability and trustworthiness in real-world applications. Furthermore, the training of large language models requires significant computational resources and large amounts of data, which can be a barrier to their adoption in resource-constrained environments.

To overcome these challenges, researchers have proposed various techniques, such as model pruning, quantization, and knowledge distillation, to reduce the computational costs and improve the efficiency of large language models. For example, \cite{arxiv:2305.16264} discusses the use of model pruning and quantization to reduce the size and computational requirements of large language models, while \cite{arxiv:2402.01364} explores the application of continual learning techniques to adapt large language models to new tasks and environments.

In addition to these technical challenges, the development and deployment of large language model-based agents also raise important ethical and societal considerations. As highlighted in \cite{doi:10.1145/3604237.3626869}, the use of large language models in financial applications can have significant implications for financial stability and fairness, underscoring the need for careful evaluation and regulation of these models. Similarly, \cite{doi:10.1109/mnet.2024.3435752} discusses the potential applications of large language models in networking and cybersecurity, emphasizing the need for robust security measures to prevent the misuse of these models.

In conclusion, the development and application of large language model-based agents represent a rapidly evolving field with significant potential for innovation and impact. While challenges and limitations remain to be addressed, the continued advancement of these models is likely to have far-reaching implications for various domains, from natural language processing and computer vision to finance and healthcare. As noted in \cite{doi:10.1109/tkde.2025.3554028}, the future of large language models is likely to be shaped by the development of more efficient and effective techniques for training and deploying these models, as well as the integration of these models with other technologies, such as computer vision and robotics. Ultimately, the successful development and deployment of large language model-based agents will require a multidisciplinary approach, combining advances in machine learning, natural language processing, and human-computer interaction with careful consideration of the ethical and societal implications of these technologies, as emphasized in \cite{arxiv:2405.03207} and \cite{arxiv:2404.01023}.

\begin{thebibliography}{109}
\setlength{\itemsep}{2pt}

\bibitem{arxiv:2206.07682}
\newblock \emph{Emergent Abilities of Large Language Models}.
\newblock arXiv:2206.07682, 2022.
\newblock \url{https://arxiv.org/abs/2206.07682}

\bibitem{doi:10.1016/j.metrad.2023.100017}
\newblock \emph{Summary of ChatGPT-Related Research and Perspective Towards the Future of Large Language Models}.
\newblock DOI:10.1016/j.metrad.2023.100017, 2023.
\newblock \url{https://doi.org/10.1016/j.metrad.2023.100017}

\bibitem{doi:10.1145/3604237.3626869}
\newblock \emph{Large Language Models in Finance: A Survey}.
\newblock DOI:10.1145/3604237.3626869, 2023.
\newblock \url{https://doi.org/10.1145/3604237.3626869}

\bibitem{arxiv:2305.16264}
\newblock \emph{Scaling Data-Constrained Language Models}.
\newblock arXiv:2305.16264, 2023.
\newblock \url{https://arxiv.org/abs/2305.16264}

\bibitem{arxiv:2309.03409}
\newblock \emph{Large Language Models as Optimizers}.
\newblock arXiv:2309.03409, 2023.
\newblock \url{https://arxiv.org/abs/2309.03409}

\bibitem{doi:10.1007/s10849-023-09409-x}
\newblock \emph{Assessing the Strengths and Weaknesses of Large Language Models}.
\newblock DOI:10.1007/s10849-023-09409-x, 2023.
\newblock \url{https://doi.org/10.1007/s10849-023-09409-x}

\bibitem{doi:10.18653/v1/2023.acl-long.411}
\newblock \emph{Large Language Models Meet NL2Code: A Survey}.
\newblock DOI:10.18653/v1/2023.acl-long.411, 2023.
\newblock \url{https://doi.org/10.18653/v1/2023.acl-long.411}

\bibitem{arxiv:2401.13601}
\newblock \emph{MM-LLMs: Recent Advances in MultiModal Large Language Models}.
\newblock arXiv:2401.13601, 2024.
\newblock \url{https://arxiv.org/abs/2401.13601}

\bibitem{doi:10.1109/mnet.2024.3435752}
\newblock \emph{Large Language Models for Networking: Applications, Enabling Techniques, and Challenges}.
\newblock DOI:10.1109/mnet.2024.3435752, 2024.
\newblock \url{https://doi.org/10.1109/mnet.2024.3435752}

\bibitem{doi:10.1109/tkde.2025.3554028}
\newblock \emph{A Survey on Mixture of Experts in Large Language Models}.
\newblock DOI:10.1109/tkde.2025.3554028, 2025.
\newblock \url{https://doi.org/10.1109/tkde.2025.3554028}

\bibitem{doi:10.1109/mwc.016.2300600}
\newblock \emph{Large Language Model Enhanced Multi-Agent Systems for 6G Communications}.
\newblock DOI:10.1109/mwc.016.2300600, 2024.
\newblock \url{https://doi.org/10.1109/mwc.016.2300600}

\bibitem{arxiv:2309.14365}
\newblock \emph{An In-depth Survey of Large Language Model-based Artificial Intelligence Agents}.
\newblock arXiv:2309.14365, 2023.
\newblock \url{https://arxiv.org/abs/2309.14365}

\bibitem{arxiv:2308.03688}
\newblock \emph{AgentBench: Evaluating LLMs as Agents}.
\newblock arXiv:2308.03688, 2023.
\newblock \url{https://arxiv.org/abs/2308.03688}

\bibitem{arxiv:2402.01364}
\newblock \emph{Continual Learning for Large Language Models: A Survey}.
\newblock arXiv:2402.01364, 2024.
\newblock \url{https://arxiv.org/abs/2402.01364}

\bibitem{doi:10.1109/access.2024.3424945}
\newblock \emph{Achieving Peak Performance for Large Language Models: A Systematic Review}.
\newblock DOI:10.1109/access.2024.3424945, 2024.
\newblock \url{https://doi.org/10.1109/access.2024.3424945}

\bibitem{doi:10.1145/3641289}
\newblock \emph{A Survey on Evaluation of Large Language Models}.
\newblock DOI:10.1145/3641289, 2024.
\newblock \url{https://doi.org/10.1145/3641289}

\bibitem{arxiv:2309.02427}
\newblock \emph{Cognitive Architectures for Language Agents}.
\newblock arXiv:2309.02427, 2023.
\newblock \url{https://arxiv.org/abs/2309.02427}

\bibitem{doi:10.3390/app14052074}
\newblock \emph{A Review of Current Trends, Techniques, and Challenges in Large Language Models (LLMs)}.
\newblock DOI:10.3390/app14052074, 2024.
\newblock \url{https://doi.org/10.3390/app14052074}

\bibitem{doi:10.34133/research.0646}
\newblock \emph{When Large Language Models Meet Evolutionary Algorithms: Potential Enhancements and Challenges}.
\newblock DOI:10.34133/research.0646, 2025.
\newblock \url{https://doi.org/10.34133/research.0646}

\bibitem{arxiv:1412.6980}
\newblock \emph{Adam: A Method for Stochastic Optimization}.
\newblock arXiv:1412.6980, 2014.
\newblock \url{https://arxiv.org/abs/1412.6980}

\bibitem{arxiv:2402.15627}
\newblock \emph{MegaScale: Scaling Large Language Model Training to More Than 10, 000 GPUs}.
\newblock arXiv:2402.15627, 2024.
\newblock \url{https://arxiv.org/abs/2402.15627}

\bibitem{doi:10.1145/3748302}
\newblock \emph{A Survey on the Memory Mechanism of Large Language Model-based Agents}.
\newblock DOI:10.1145/3748302, 2025.
\newblock \url{https://doi.org/10.1145/3748302}

\bibitem{arxiv:2401.02777}
\newblock \emph{From LLM to Conversational Agent: A Memory Enhanced Architecture with Fine-Tuning of Large Language Models}.
\newblock arXiv:2401.02777, 2024.
\newblock \url{https://arxiv.org/abs/2401.02777}

\bibitem{arxiv:2410.02644}
\newblock \emph{Agent Security Bench (ASB): Formalizing and Benchmarking Attacks and Defenses in LLM-based Agents}.
\newblock arXiv:2410.02644, 2024.
\newblock \url{https://arxiv.org/abs/2410.02644}

\bibitem{arxiv:1906.09379}
\newblock \emph{Evaluating Computational Language Models with Scaling Properties of Natural Language}.
\newblock arXiv:1906.09379, 2019.
\newblock \url{https://arxiv.org/abs/1906.09379}

\bibitem{arxiv:2308.04026}
\newblock \emph{AgentSims: An Open-Source Sandbox for Large Language Model Evaluation}.
\newblock arXiv:2308.04026, 2023.
\newblock \url{https://arxiv.org/abs/2308.04026}

\bibitem{doi:10.18653/v1/2023.findings-emnlp.128}
\newblock \emph{Transformer-Based Language Model Surprisal Predicts Human Reading Times Best with About Two Billion Training Tokens}.
\newblock DOI:10.18653/v1/2023.findings-emnlp.128, 2023.
\newblock \url{https://doi.org/10.18653/v1/2023.findings-emnlp.128}

\bibitem{arxiv:2301.04589}
\newblock \emph{Memory Augmented Large Language Models are Computationally Universal}.
\newblock arXiv:2301.04589, 2023.
\newblock \url{https://arxiv.org/abs/2301.04589}

\bibitem{arxiv:2502.10708}
\newblock \emph{Injecting Domain-Specific Knowledge into Large Language Models: A Comprehensive Survey}.
\newblock arXiv:2502.10708, 2025.
\newblock \url{https://arxiv.org/abs/2502.10708}

\bibitem{arxiv:2402.12914}
\newblock \emph{Large Language Model-based Human-Agent Collaboration for Complex Task Solving}.
\newblock arXiv:2402.12914, 2024.
\newblock \url{https://arxiv.org/abs/2402.12914}

\bibitem{arxiv:2402.11208}
\newblock \emph{Watch Out for Your Agents! Investigating Backdoor Threats to LLM-Based Agents}.
\newblock arXiv:2402.11208, 2024.
\newblock \url{https://arxiv.org/abs/2402.11208}

\bibitem{arxiv:2406.04151}
\newblock \emph{AgentGym: Evolving Large Language Model-based Agents across Diverse Environments}.
\newblock arXiv:2406.04151, 2024.
\newblock \url{https://arxiv.org/abs/2406.04151}

\bibitem{doi:10.22214/ijraset.2023.54677}
\newblock \emph{Several categories of Large Language Models (LLMs): A Short Survey}.
\newblock DOI:10.22214/ijraset.2023.54677, 2023.
\newblock \url{https://doi.org/10.22214/ijraset.2023.54677}

\bibitem{arxiv:2308.05960}
\newblock \emph{BOLAA: Benchmarking and Orchestrating LLM-augmented Autonomous Agents}.
\newblock arXiv:2308.05960, 2023.
\newblock \url{https://arxiv.org/abs/2308.05960}

\bibitem{arxiv:2308.02151}
\newblock \emph{Retroformer: Retrospective Large Language Agents with Policy Gradient Optimization}.
\newblock arXiv:2308.02151, 2023.
\newblock \url{https://arxiv.org/abs/2308.02151}

\bibitem{arxiv:2311.13884}
\newblock \emph{Controlling Large Language Model-based Agents for Large-Scale Decision-Making: An Actor-Critic Approach}.
\newblock arXiv:2311.13884, 2023.
\newblock \url{https://arxiv.org/abs/2311.13884}

\bibitem{arxiv:2412.14470}
\newblock \emph{Agent-SafetyBench: Evaluating the Safety of LLM Agents}.
\newblock arXiv:2412.14470, 2024.
\newblock \url{https://arxiv.org/abs/2412.14470}

\bibitem{arxiv:2501.05707}
\newblock \emph{Multiagent Finetuning: Self Improvement with Diverse Reasoning Chains}.
\newblock arXiv:2501.05707, 2025.
\newblock \url{https://arxiv.org/abs/2501.05707}

\bibitem{arxiv:2402.16823}
\newblock \emph{Language Agents as Optimizable Graphs}.
\newblock arXiv:2402.16823, 2024.
\newblock \url{https://arxiv.org/abs/2402.16823}

\bibitem{arxiv:2406.04692}
\newblock \emph{Mixture-of-Agents Enhances Large Language Model Capabilities}.
\newblock arXiv:2406.04692, 2024.
\newblock \url{https://arxiv.org/abs/2406.04692}

\bibitem{doi:10.1609/aaaiss.v2i1.27691}
\newblock \emph{A Proposal for a Language Model Based Cognitive Architecture}.
\newblock DOI:10.1609/aaaiss.v2i1.27691, 2024.
\newblock \url{https://doi.org/10.1609/aaaiss.v2i1.27691}

\bibitem{arxiv:2402.01680}
\newblock \emph{Large Language Model based Multi-Agents: A Survey of Progress and Challenges}.
\newblock arXiv:2402.01680, 2024.
\newblock \url{https://arxiv.org/abs/2402.01680}

\bibitem{arxiv:2403.14608}
\newblock \emph{Parameter-Efficient Fine-Tuning for Large Models: A Comprehensive Survey}.
\newblock arXiv:2403.14608, 2024.
\newblock \url{https://arxiv.org/abs/2403.14608}

\bibitem{arxiv:2402.06196}
\newblock \emph{Large Language Models: A Survey}.
\newblock arXiv:2402.06196, 2024.
\newblock \url{https://arxiv.org/abs/2402.06196}

\bibitem{doi:10.1109/iccv51070.2023.00788}
\newblock \emph{MotionLM: Multi-Agent Motion Forecasting as Language Modeling}.
\newblock DOI:10.1109/iccv51070.2023.00788, 2023.
\newblock \url{https://doi.org/10.1109/iccv51070.2023.00788}

\bibitem{arxiv:2312.00678}
\newblock \emph{The Efficiency Spectrum of Large Language Models: An Algorithmic Survey}.
\newblock arXiv:2312.00678, 2023.
\newblock \url{https://arxiv.org/abs/2312.00678}

\bibitem{doi:10.3348/kjr.2023.0997}
\newblock \emph{Large Language Models: A Guide for Radiologists}.
\newblock DOI:10.3348/kjr.2023.0997, 2024.
\newblock \url{https://doi.org/10.3348/kjr.2023.0997}

\bibitem{doi:10.1016/j.tics.2023.08.006}
\newblock \emph{What are large language models supposed to model?}.
\newblock DOI:10.1016/j.tics.2023.08.006, 2023.
\newblock \url{https://doi.org/10.1016/j.tics.2023.08.006}

\bibitem{arxiv:2310.20151}
\newblock \emph{Multi-Agent Consensus Seeking via Large Language Models}.
\newblock arXiv:2310.20151, 2023.
\newblock \url{https://arxiv.org/abs/2310.20151}

\bibitem{arxiv:2410.13825}
\newblock \emph{AgentOccam: A Simple Yet Strong Baseline for LLM-Based Web Agents}.
\newblock arXiv:2410.13825, 2024.
\newblock \url{https://arxiv.org/abs/2410.13825}

\bibitem{doi:10.1109/mnet.2025.3539338}
\newblock \emph{Toward Edge General Intelligence via Large Language Models: Opportunities and Challenges}.
\newblock DOI:10.1109/mnet.2025.3539338, 2025.
\newblock \url{https://doi.org/10.1109/mnet.2025.3539338}

\bibitem{doi:10.1145/3649506}
\newblock \emph{Harnessing the Power of LLMs in Practice: A Survey on ChatGPT and Beyond}.
\newblock DOI:10.1145/3649506, 2024.
\newblock \url{https://doi.org/10.1145/3649506}

\bibitem{arxiv:2402.10172}
\newblock \emph{OptiMUS: Scalable Optimization Modeling with (MI)LP Solvers and Large Language Models}.
\newblock arXiv:2402.10172, 2024.
\newblock \url{https://arxiv.org/abs/2402.10172}

\bibitem{arxiv:2409.06857}
\newblock \emph{What is the Role of Small Models in the LLM Era: A Survey}.
\newblock arXiv:2409.06857, 2024.
\newblock \url{https://arxiv.org/abs/2409.06857}

\bibitem{arxiv:2402.05120}
\newblock \emph{More Agents Is All You Need}.
\newblock arXiv:2402.05120, 2024.
\newblock \url{https://arxiv.org/abs/2402.05120}

\bibitem{doi:10.34133/research.0655}
\newblock \emph{Road of Large Language Model: Source, Challenge, and Future Perspectives}.
\newblock DOI:10.34133/research.0655, 2025.
\newblock \url{https://doi.org/10.34133/research.0655}

\bibitem{arxiv:2404.01023}
\newblock \emph{Large Language Model Evaluation Via Multi AI Agents: Preliminary results}.
\newblock arXiv:2404.01023, 2024.
\newblock \url{https://arxiv.org/abs/2404.01023}

\bibitem{doi:10.1162/coli_a_00107}
\newblock \emph{A Scalable Distributed Syntactic, Semantic, and Lexical Language Model}.
\newblock DOI:10.1162/coli\_a\_00107, 2012.
\newblock \url{https://doi.org/10.1162/coli_a_00107}

\bibitem{arxiv:2112.11446}
\newblock \emph{Scaling Language Models: Methods, Analysis \& Insights from Training Gopher}.
\newblock arXiv:2112.11446, 2021.
\newblock \url{https://arxiv.org/abs/2112.11446}

\bibitem{doi:10.18653/v1/2022.bigscience-1.1}
\newblock \emph{Lifelong Pretraining: Continually Adapting Language Models to Emerging Corpora}.
\newblock DOI:10.18653/v1/2022.bigscience-1.1, 2022.
\newblock \url{https://doi.org/10.18653/v1/2022.bigscience-1.1}

\bibitem{doi:10.18653/v1/2021.naacl-main.341}
\newblock \emph{Progressive Generation of Long Text with Pretrained Language Models}.
\newblock DOI:10.18653/v1/2021.naacl-main.341, 2021.
\newblock \url{https://doi.org/10.18653/v1/2021.naacl-main.341}

\bibitem{arxiv:2304.05332}
\newblock \emph{Emergent autonomous scientific research capabilities of large language models}.
\newblock arXiv:2304.05332, 2023.
\newblock \url{https://arxiv.org/abs/2304.05332}

\bibitem{arxiv:2205.05718}
\newblock \emph{Structured, flexible, and robust: benchmarking and improving large language models towards more human-like behavior in out-of-distribution reasoning tasks}.
\newblock arXiv:2205.05718, 2022.
\newblock \url{https://arxiv.org/abs/2205.05718}

\bibitem{arxiv:2306.03917}
\newblock \emph{Turning large language models into cognitive models}.
\newblock arXiv:2306.03917, 2023.
\newblock \url{https://arxiv.org/abs/2306.03917}

\bibitem{arxiv:2407.07086}
\newblock \emph{Hypothetical Minds: Scaffolding Theory of Mind for Multi-Agent Tasks with Large Language Models}.
\newblock arXiv:2407.07086, 2024.
\newblock \url{https://arxiv.org/abs/2407.07086}

\bibitem{arxiv:2303.18223}
\newblock \emph{A Survey of Large Language Models}.
\newblock arXiv:2303.18223, 2023.
\newblock \url{https://arxiv.org/abs/2303.18223}

\bibitem{arxiv:2311.06622}
\newblock \emph{TrainerAgent: Customizable and Efficient Model Training through LLM-Powered Multi-Agent System}.
\newblock arXiv:2311.06622, 2023.
\newblock \url{https://arxiv.org/abs/2311.06622}

\bibitem{doi:10.1609/aaaiss.v2i1.27688}
\newblock \emph{Memory Matters: The Need to Improve Long-Term Memory in LLM-Agents}.
\newblock DOI:10.1609/aaaiss.v2i1.27688, 2024.
\newblock \url{https://doi.org/10.1609/aaaiss.v2i1.27688}

\bibitem{arxiv:2401.06201}
\newblock \emph{EASYTOOL: Enhancing LLM-based Agents with Concise Tool Instruction}.
\newblock arXiv:2401.06201, 2024.
\newblock \url{https://arxiv.org/abs/2401.06201}

\bibitem{arxiv:2410.06153}
\newblock \emph{AgentSquare: Automatic LLM Agent Search in Modular Design Space}.
\newblock arXiv:2410.06153, 2024.
\newblock \url{https://arxiv.org/abs/2410.06153}

\bibitem{arxiv:2401.10910}
\newblock \emph{Metacognition is all you need? Using Introspection in Generative Agents to Improve Goal-directed Behavior}.
\newblock arXiv:2401.10910, 2024.
\newblock \url{https://arxiv.org/abs/2401.10910}

\bibitem{arxiv:2404.09135}
\newblock \emph{Unveiling LLM Evaluation Focused on Metrics: Challenges and Solutions}.
\newblock arXiv:2404.09135, 2024.
\newblock \url{https://arxiv.org/abs/2404.09135}

\bibitem{arxiv:2404.09022}
\newblock \emph{Navigating the Landscape of Large Language Models: A Comprehensive Review and Analysis of Paradigms and Fine-Tuning Strategies}.
\newblock arXiv:2404.09022, 2024.
\newblock \url{https://arxiv.org/abs/2404.09022}

\bibitem{doi:10.1145/3637528.3671650}
\newblock \emph{Understanding the Weakness of Large Language Model Agents within a Complex Android Environment}.
\newblock DOI:10.1145/3637528.3671650, 2024.
\newblock \url{https://doi.org/10.1145/3637528.3671650}

\bibitem{doi:10.1109/bigdata59044.2023.10386743}
\newblock \emph{Multimodal Large Language Models: A Survey}.
\newblock DOI:10.1109/bigdata59044.2023.10386743, 2023.
\newblock \url{https://doi.org/10.1109/bigdata59044.2023.10386743}

\bibitem{doi:10.1109/tnnls.2024.3497992}
\newblock \emph{Survey on Large Language Model-Enhanced Reinforcement Learning: Concept, Taxonomy, and Methods}.
\newblock DOI:10.1109/tnnls.2024.3497992, 2024.
\newblock \url{https://doi.org/10.1109/tnnls.2024.3497992}

\bibitem{doi:10.1145/3754448}
\newblock \emph{Towards Efficient Generative Large Language Model Serving: A Survey from Algorithms to Systems}.
\newblock DOI:10.1145/3754448, 2025.
\newblock \url{https://doi.org/10.1145/3754448}

\bibitem{doi:10.1109/access.2024.3365742}
\newblock \emph{A Review on Large Language Models: Architectures, Applications, Taxonomies, Open Issues and Challenges}.
\newblock DOI:10.1109/access.2024.3365742, 2024.
\newblock \url{https://doi.org/10.1109/access.2024.3365742}

\bibitem{arxiv:2212.09251}
\newblock \emph{Discovering Language Model Behaviors with Model-Written Evaluations}.
\newblock arXiv:2212.09251, 2022.
\newblock \url{https://arxiv.org/abs/2212.09251}

\bibitem{arxiv:2402.14744}
\newblock \emph{Large Language Models as Urban Residents: An LLM Agent Framework for Personal Mobility Generation}.
\newblock arXiv:2402.14744, 2024.
\newblock \url{https://arxiv.org/abs/2402.14744}

\bibitem{arxiv:2402.16142}
\newblock \emph{From Text to Transformation: A Comprehensive Review of Large Language Models' Versatility}.
\newblock arXiv:2402.16142, 2024.
\newblock \url{https://arxiv.org/abs/2402.16142}

\bibitem{doi:10.1038/s41746-024-01083-y}
\newblock \emph{Evaluating large language models as agents in the clinic}.
\newblock DOI:10.1038/s41746-024-01083-y, 2024.
\newblock \url{https://doi.org/10.1038/s41746-024-01083-y}

\bibitem{arxiv:2405.11357}
\newblock \emph{Large Language Models Lack Understanding of Character Composition of Words}.
\newblock arXiv:2405.11357, 2024.
\newblock \url{https://arxiv.org/abs/2405.11357}

\bibitem{arxiv:2404.06411}
\newblock \emph{AgentQuest: A Modular Benchmark Framework to Measure Progress and Improve LLM Agents}.
\newblock arXiv:2404.06411, 2024.
\newblock \url{https://arxiv.org/abs/2404.06411}

\bibitem{doi:10.1007/s11704-024-40231-1}
\newblock \emph{A survey on large language model based autonomous agents}.
\newblock DOI:10.1007/s11704-024-40231-1, 2024.
\newblock \url{https://doi.org/10.1007/s11704-024-40231-1}

\bibitem{arxiv:2411.04890}
\newblock \emph{GUI Agents with Foundation Models: A Comprehensive Survey}.
\newblock arXiv:2411.04890, 2024.
\newblock \url{https://arxiv.org/abs/2411.04890}

\bibitem{arxiv:2501.11354}
\newblock \emph{Towards Advancing Code Generation with Large Language Models: A Research Roadmap}.
\newblock arXiv:2501.11354, 2025.
\newblock \url{https://arxiv.org/abs/2501.11354}

\bibitem{arxiv:2402.11443}
\newblock \emph{Benchmark Self-Evolving: A Multi-Agent Framework for Dynamic LLM Evaluation}.
\newblock arXiv:2402.11443, 2024.
\newblock \url{https://arxiv.org/abs/2402.11443}

\bibitem{doi:10.1145/3624724}
\newblock \emph{Talking about Large Language Models}.
\newblock DOI:10.1145/3624724, 2024.
\newblock \url{https://doi.org/10.1145/3624724}

\bibitem{doi:10.1007/s13042-024-02318-w}
\newblock \emph{Large language models for medicine: a survey}.
\newblock DOI:10.1007/s13042-024-02318-w, 2024.
\newblock \url{https://doi.org/10.1007/s13042-024-02318-w}

\bibitem{doi:10.1145/3605943}
\newblock \emph{Recent Advances in Natural Language Processing via Large Pre-trained Language Models: A Survey}.
\newblock DOI:10.1145/3605943, 2023.
\newblock \url{https://doi.org/10.1145/3605943}

\bibitem{arxiv:2403.19887}
\newblock \emph{Jamba: A Hybrid Transformer-Mamba Language Model}.
\newblock arXiv:2403.19887, 2024.
\newblock \url{https://arxiv.org/abs/2403.19887}

\bibitem{arxiv:2408.06361}
\newblock \emph{Large Language Model Agent in Financial Trading: A Survey}.
\newblock arXiv:2408.06361, 2024.
\newblock \url{https://arxiv.org/abs/2408.06361}

\bibitem{arxiv:2305.17144}
\newblock \emph{Ghost in the Minecraft: Generally Capable Agents for Open-World Environments via Large Language Models with Text-based Knowledge and Memory}.
\newblock arXiv:2305.17144, 2023.
\newblock \url{https://arxiv.org/abs/2305.17144}

\bibitem{arxiv:2309.07870}
\newblock \emph{Agents: An Open-source Framework for Autonomous Language Agents}.
\newblock arXiv:2309.07870, 2023.
\newblock \url{https://arxiv.org/abs/2309.07870}

\bibitem{doi:10.1145/3712003}
\newblock \emph{LLM-Based Multi-Agent Systems for Software Engineering: Literature Review, Vision, and the Road Ahead}.
\newblock DOI:10.1145/3712003, 2025.
\newblock \url{https://doi.org/10.1145/3712003}

\bibitem{arxiv:2402.01108}
\newblock \emph{Reasoning Capacity in Multi-Agent Systems: Limitations, Challenges and Human-Centered Solutions}.
\newblock arXiv:2402.01108, 2024.
\newblock \url{https://arxiv.org/abs/2402.01108}

\bibitem{arxiv:2402.00798}
\newblock \emph{Formal-LLM: Integrating Formal Language and Natural Language for Controllable LLM-based Agents}.
\newblock arXiv:2402.00798, 2024.
\newblock \url{https://arxiv.org/abs/2402.00798}

\bibitem{arxiv:2402.01881}
\newblock \emph{Large Language Model Agent for Hyper-Parameter Optimization}.
\newblock arXiv:2402.01881, 2024.
\newblock \url{https://arxiv.org/abs/2402.01881}

\bibitem{doi:10.1109/access.2024.3524176}
\newblock \emph{Metaheuristics and Large Language Models Join Forces: Toward an Integrated Optimization Approach}.
\newblock DOI:10.1109/access.2024.3524176, 2024.
\newblock \url{https://doi.org/10.1109/access.2024.3524176}

\bibitem{arxiv:2405.19616}
\newblock \emph{Easy Problems That LLMs Get Wrong}.
\newblock arXiv:2405.19616, 2024.
\newblock \url{https://arxiv.org/abs/2405.19616}

\bibitem{doi:10.1109/mwc.005.2400019}
\newblock \emph{When Large Language Model Agents Meet 6G Networks: Perception, Grounding, and Alignment}.
\newblock DOI:10.1109/mwc.005.2400019, 2024.
\newblock \url{https://doi.org/10.1109/mwc.005.2400019}

\bibitem{arxiv:2405.03207}
\newblock \emph{A Philosophical Introduction to Language Models - Part II: The Way Forward}.
\newblock arXiv:2405.03207, 2024.
\newblock \url{https://arxiv.org/abs/2405.03207}

\bibitem{arxiv:2311.07226}
\newblock \emph{Large Language Models for Robotics: A Survey}.
\newblock arXiv:2311.07226, 2023.
\newblock \url{https://arxiv.org/abs/2311.07226}

\bibitem{doi:10.54254/2755-2721/35/20230399}
\newblock \emph{The evolution, applications, and future prospects of large language models: An in-depth overview}.
\newblock DOI:10.54254/2755-2721/35/20230399, 2024.
\newblock \url{https://doi.org/10.54254/2755-2721/35/20230399}

\bibitem{arxiv:2311.07621}
\newblock \emph{To Transformers and Beyond: Large Language Models for the Genome}.
\newblock arXiv:2311.07621, 2023.
\newblock \url{https://arxiv.org/abs/2311.07621}

\bibitem{doi:10.18653/v1/2023.findings-emnlp.461}
\newblock \emph{Towards large language model-based personal agents in the enterprise: Current trends and open problems}.
\newblock DOI:10.18653/v1/2023.findings-emnlp.461, 2023.
\newblock \url{https://doi.org/10.18653/v1/2023.findings-emnlp.461}

\bibitem{arxiv:2308.04030}
\newblock \emph{Gentopia: A Collaborative Platform for Tool-Augmented LLMs}.
\newblock arXiv:2308.04030, 2023.
\newblock \url{https://arxiv.org/abs/2308.04030}

\bibitem{arxiv:2408.15769}
\newblock \emph{A Survey on Evaluation of Multimodal Large Language Models}.
\newblock arXiv:2408.15769, 2024.
\newblock \url{https://arxiv.org/abs/2408.15769}

\end{thebibliography}

\end{document}
